% =====================================================================
% verna.tex -- LaTeX conversion of verna.md (Breviarium Romanum,
% Pars Verna), BODY CONTENT ONLY.
%
% No preamble, \documentclass, or \newcommand definitions are included
% here, per instruction -- the command vocabulary below is to be
% defined locally, following the same legend established for
% hiemalis.tex (one-column vs. two-column, red vs. black, etc.).
%
% BATCH 1 of N: Ordinarium Divini Officii (verna.md lines 1-1745)
% =====================================================================

% [ornamental image beginning Pars Verna]
\partheading{Ordinarium Divini Officii}
\partsubtitle{Juxta Ritum Romanum Persolvendi}

\rubric{Clericis in sacris constitútis, qui íntegrum divínum Offícium, quamvis in partes distribútum, coram Ssmo Sacraménto, sive públicæ adoratióni expósito sive in tabernáculo adserváto, devóte recitáverint, concéditur: Indulgéntia plenária, si prætérea sacramentálem confessiónem instítuerint, ad eucharísticum Convívium accésserint et ad mentem Summi Pontíficis oráverint. \textit{Pius Pp. XI, 23 Octóbris 1930.}}

\rubric{Iis vero, qui divínum Offícium coram Ssmo Sacraménto, ut supra, ex parte tantum recitáverint, concéditur: Indulgéntia quingentórum diérum pro unaquáque hora canónica. \textit{Pius Pp. XI, 18 Maji 1933.}}

\rubric{Iisdem autem cléricis, qui divínum Offícium in álias preces commutátum rite obtinúerint, si preces ejúsmodi coram Ssmo Sacraménto, uti supra, devóte recitáverint, concéditur: Indulgéntia plenária, addítis sacramentáli confessióne, sacra Communióne et oratióne ad mentem Summi Pontíficis. \textit{Pius Pp. XI, 7 Novémbris 1932.}}

\sectionheading{Ante Divinum Officium}

\rubric{Antequam inchoëtur Offícium laudabíliter dícitur, sub singulári semper número, séquens Orátio; pro qua Summus Póntifex Pius Pp. XI sub die 17 Novémbris 1933 Indulgéntiam trium annórum concéssit.}

\oratioheading{Oratio}

\lettrine[lines=2]{A}{PERI}, Dómine, os meum ad benedicéndum nomen sanctum tuum: munda quoque cor meum ab ómnibus vanis, pervérsis et aliénis cogitatiónibus; intelléctum illúmina, afféctum inflámma, ut digne, atténte ac devóte hoc Offícium recitáre váleam, et exaudíri mérear ante conspéctum divínæ Majestátis tuæ. Per Christum, Dóminum nostrum. \R\ Amen.

Dómine, in unióne illíus divínæ intentiónis, qua ipse in terris laudes Deo persolvísti, has tibi Horas (\textit{vel} hanc tibi Horam) persólvo.

\sectionheading{Ante singulas Horas}

\rubric{Ante Matutínum, sicut et ante Laudes a Matutíno extra Chorum separátas, et ante omnes Horas, præter quam ad Completórium, dícitur secréto:}

\lettrine[lines=2]{P}{ATER} noster, qui es in cælis, sanctificétur nomen tuum. Advéniat regnum tuum. Fiat volúntas tua, sicut in cælo et in terra. Panem nostrum quotidiánum da nobis hódie. Et dimítte nobis débita nostra, sicut et nos dimíttimus debitóribus nostris. Et ne nos indúcas in tentatiónem: sed líbera nos a malo. Amen.

\lettrine[lines=2]{A}{VE}, María, grátia plena; Dóminus tecum: benedícta tu in muliéribus, et benedíctus fructus ventris tui Jesus. Sancta María, Mater Dei, ora pro nobis peccatóribus, nunc et in hora mortis nostræ. Amen.

\rubric{In princípio autem Matutíni ac Primæ additur secréto Symbolum Apostolórum:}

\lettrine[lines=2]{C}{REDO} in Deum, Patrem omnipoténtem, Creatórem cæli et terræ. Et in Jesum Christum Fílium ejus únicum, Dóminum nostrum: qui concéptus est de Spíritu Sancto, natus ex María Vírgine, passus sub Póntio Piláto, crucifíxus, mórtuus, et sepúltus: descéndit ad ínferos; tértia die resurréxit a mórtuis; ascéndit ad cælos; sedet ad déxteram Dei Patris omnipoténtis: inde ventúrus est judicáre vivos et mórtuos. Credo in Spíritum Sanctum, sanctam Ecclésiam cathólicam, Sanctórum communiónem, remissiónem peccatórum, carnis resurrectiónem, vitam ætérnam. Amen.

\rubric{Offícium ad Matutínum et ad réliquas Horas, nisi áliter suo loco indicétur, per totum annum persólvitur juxta formam, quæ in Rúbricis hujus Ordinárii exhibétur.}

\hourheading{Ad Matutinum}

Pater noster, Ave María et Credo.

\rubric{Deinde, clara voce, dicitur Versus:}

\V\ \lettrine[lines=2]{D}{OMINE}, lábia mea apéries. \R\ Et os meum annuntiábit laudem tuam. \V\ Deus, in adjutórium meum inténde. \R\ Dómine, ad adjuvándum me festína.

Glória Patri, et Fílio, et Spirítui Sancto. Sicut erat in princípio, et nunc, et semper, et in sæcula sæculórum. Amen. Allelúja.

\rubric{Sic dicitur Allelúja ad omnes Horas, præter quam a Vésperis Sábbati ante Dominícam I Quadragésimæ usque ad Completórium Fériæ IV Majóris Hebdómadæ inclusíve; tunc enim dicitur:}

Laus tibi, Dómine, Rex ætérnæ glóriæ.

\rubric{Postea dicitur convéniens Invitatórium, quod ante Psalmum bis íntegrum recitátur, et ad síngulos ejúsdem Psalmi versus vel íntegrum vel dimidiátum ab Asterísco * altérnis vícibus, ut infra, repétitur.}

\psalmheading{Psalmus 94}

\lettrine[lines=2]{V}{ENITE}, exsultémus Dómino, jubilémus Deo, salutári nostro: præoccupémus fáciem ejus in confessióne, et in psalmis jubilémus ei.

\rubric{Repetitur integrum Invitatorium.}

Quóniam Deus magnus Dóminus, et Rex magnus super omnes deos: quóniam non repéllet Dóminus plebem suam: quia in manu ejus sunt omnes fines terræ, et altitúdines móntium ipse cónspicit.

\rubric{Repetitur altera pars Invitatorii.}

\rubric{In sequenti Psalmi versu, ad verba veníte, adorémus, et procidámus ante Deum, genuflectitur.}

Quóniam ipsíus est mare, et ipse fecit illud, et áridam fundavérunt manus ejus: veníte, adorémus, et procidámus ante Deum: plorémus coram Dómino, qui fecit nos, quia ipse est Dóminus, Deus noster; nos autem pópulus ejus, et oves páscuæ ejus.

\rubric{Repetitur integrum Invitatorium.}

Hódie, si vocem ejus audíeritis, nolíte obduráre corda vestra, sicut in exacerbatióne secúndum diem tentatiónis in desérto: ubi tentavérunt me patres vestri, probavérunt et vidérunt ópera mea.

\rubric{Repetitur altera pars Invitatorii.}

Quadragínta annis próximus fui generatióni huic, et dixi: Semper hi errant corde; ipsi vero non cognovérunt vias meas: quibus jurávi in ira mea: Si introíbunt in réquiem meam.

\rubric{Repetitur integrum Invitatorium.}

Glória Patri, et Fílio, et Spirítui Sancto. Sicut erat in princípio, et nunc, et semper, et in sæcula sæculórum. Amen.

\rubric{Repetitur altera pars Invitatorii.}

\rubric{Denique repetitur integrum Invitatorium.}

\rubric{Expleto Psalmo, dicitur Hymnus Invitatorio respondens.}

\variantheading{Tempore Quadragesimæ}

\rubric{A Dominica I usque ad Sabbatum ante Dominicam Passiónis inclusíve, tam in dominicáli quam in feriáli Officio, dícitur quotídie séquens}

\invitatlabel{Invitat.} Non sit vobis vanum mane súrgere ante lucem: * Quia promísit Dóminus corónam vigilántibus. \psalmlabel{Ps. 94} Veníte. [3]

\hymnusheading{Hymnus}

\lettrine[lines=2]{E}{X} more docti mýstico\\
Servémus hoc jejúnium,\\
Deno diérum círculo\\
Ducto quater notíssimo.

Lex et prophétæ prímitus\\
Hoc prætulérunt, póstmodum\\
Christus sacrávit, ómnium\\
Rex atque factor témporum.

Utámur ergo párcius\\
Verbis, cibis et pótibus,\\
Somno, jocis, et árctius\\
Perstémus in custódia.

Vitémus autem nóxia,\\
Quæ súbruunt mentes vagas:\\
Nullúmque demus cállidi\\
Hostis locum tyránnidi.

Flectámus iram víndicem,\\
Plorémus ante Júdicem,\\
Clamémus ore súpplici,\\
Dicámus omnes cérnui:

Nostris malis offéndimus\\
Tuam, Deus, cleméntiam:\\
Effúnde nobis désuper,\\
Remíssor, indulgéntiam.

Meménto quod sumus tui,\\
Licet cadúci, plásmatis:\\
Ne des honórem nóminis\\
Tui, precámur, álteri.

Laxa malum, quod fécimus,\\
Auge bonum, quod póscimus:\\
Placére quo tandem tibi\\
Possímus hic, et pérpetim.

\rubric{¶ Sequens Conclusio numquam mutatur.}

Præsta, beáta Trínitas,\\
Concéde, símplex Unitas,\\
Ut fructuósa sint tuis\\
Jejuniórum múnera. Amen.

\variantheading{Tempore Passionis}

\rubric{A Dominica Passiónis usque ad Fériam IV Majóris Hebdómadæ inclusíve, tam in dominicáli quam in feriáli Officio, dícitur quotídie séquens}

\invitatlabel{Invitat.} Hódie, si vocem Dómini audiéritis, * Nolíte obduráre corda vestra. \psalmlabel{Ps. 94} Veníte. [3]

\rubric{In quarto Psalmi versu non dicitur: Hódie, si vocem ejus audiéritis, nolíte obduráre corda vestra; sed, post repetitum a Choro Invitatorium, statim subjúngitur: Sicut in exacerbatióne.}

\rubric{In fine autem hujus Psalmi non dicitur Glória Patri toto Passiónis Témpore, quando fit Offícium Domínicæ aut Fériæ; sed repetitur Invitatórium, ut alias, in fine. Quod étiam servátur in Responsóriis ad Matutínum et in Responsóriis brévibus ad Horas minóres.}

\hymnusheading{Hymnus}

\lettrine[lines=2]{P}{ANGE}, lingua, gloriósi\\
Láuream certáminis,\\
Et super Crucis trophǽo\\
Dic triúmphum nóbilem:\\
Quáliter Redémptor orbis\\
Immolátus vícerit.

De paréntis protoplásti\\
Fraude Factor condolens,\\
Quando pomi noxiális\\
In necem morsu ruit:\\
Ipse lignum tunc notávit,\\
Damna ligni ut sólveret.

Hoc opus nostræ salútis\\
Ordo depopóscerat;\\
Multifórmis proditóris\\
Ars ut artem fálleret,\\
Et medélam ferret inde,\\
Hostis unde læserat.

Quando venit ergo sacri\\
Plenitúdo témporis,\\
Missus est ab arce Patris\\
Natus, orbis Cónditor;\\
Atque ventre virgináli\\
Carne amíctus pródiit.

Vagit infans inter arcta\\
Cónditus præsépia:\\
Membra pannis involúta\\
Virgo Mater álligat:\\
Et Dei manus pedésque\\
Stricta cingit fáscia.

Sempitérna sit beátæ\\
Trinitáti glória,\\
Æqua Patri, Filióque;\\
Par decus Paráclito:\\
Uníus Triníque nomen\\
Laudet univérsitas. Amen.

\variantheading{Tempore Paschali}

\rubric{A Dominica in Albis usque ad Vigíliam Ascensiónis inclusíve, tam in dominicáli quam in feriáli Officio, dícitur quotídie séquens}

\invitatlabel{Invitat.} Surréxit Dóminus vere, * Allelúja. \psalmlabel{Ps. 94} Veníte. [3]

\hymnusheading{Hymnus}

\lettrine[lines=2]{R}{EX} sempitérne cælitum,\\
Rerum Creátor ómnium,\\
Æquális ante sæcula\\
Semper Parénti Fílius:

Nascénte qui mundo faber\\
Imáginem vultus tui\\
Tradens Adámo, nóbilem\\
Limo jugásti spíritum.

Cum livor et fraus dæmonis\\
Fœdásset humánum genus:\\
Tu, carne amíctus, pérditam\\
Formam refórmas ártifex.

Qui, natus olim e Vírgine,\\
Nunc e sepúlcro násceris,\\
Tecúmque nos a mórtuis\\
Jubes sepúltos súrgere.

Qui, pastor ætérnus, gregem\\
Aqua lavas baptísmatis:\\
Hæc est lavácrum méntium:\\
Hæc est sepúlcrum críminum.

Nobis diu qui débitæ\\
Redémptor affíxus cruci,\\
Nostræ dedísti pródigus\\
Prétium salútis sánguinem.

Ut sis perénne méntibus\\
Paschále, Jesu, gáudium,\\
A morte dira críminum\\
Vitæ renátos líbera.

Deo Patri sit glória,\\
Et Fílio, qui a mórtuis\\
Surréxit, ac Paráclito,\\
In sempitérna sæcula. Amen.

\rubric{¶ Sic terminantur omnes Hymni ejúsdem metri usque ad Vigíliam Ascensiónis inclusíve, étiam in Officiis Sanctórum, nisi Conclúsio magis própria sit in eis adhibénda.}

\variantheading{In Festis}

\rubric{In Officio cujuslibet Festi vel Octávæ, ac sanctæ Maríæ in Sábbato, Invitatórium et Hymnus ut in Proprio vel Commúni; in Domínica autem infra Octávam Ascensiónis et in Féria VI post Octávam Ascensiónis, si de eis fiat Offícium, atque in Vigília Pentecóstes, Invitatórium et Hymnus ut in Proprio de Témpore. Quæ norma ad réliquas item Horas erit observánda pro illis Offícii pártibus, quæ juxta Rúbricas ex Ordinário vel Psaltério non fúerint in Festis desuméndæ.}

\variantheading{In Officio novem Lectionum}

\rubric{Expleto Hymno, dicuntur Antiphonæ convenientes, quæ in Officiis ritus duplicis ante et post Psalmos integræ recitantur; in Officiis autem ritus semiduplicis initio Psalmi inchoantur tantum et usque ad Asteriscum * perducuntur, atque in fine integræ pronuntiantur; servata alterutra norma etiam in omnibus Antiphonis, quæ ad Laudes ac Vesperas adhibentur.}

\nocturnoheading{In I Nocturno}

\rubric{Sub congruentibus Antiphonis dicuntur tres Psalmi, ac deinde subjungitur Versus, prouti Officium occurrens requirit.}

\rubric{Post Versum cujuslibet Nocturni dicitur}

Pater noster \rubric{secreto usque ad} \V\ Et ne nos indúcas in tentatiónem. \R\ Sed líbera nos a malo.

\absolutioheading{Absolutio}

\lettrine[lines=2]{E}{XÁUDI}, Dómine Jesu Christe, preces servórum tuórum, et miserére nobis: Qui cum Patre et Spíritu Sancto vivis et regnas in sæcula sæculórum. \R\ Amen.

\V\ Jube, domne, benedícere.

\rubric{¶ Extra Chorum, quando ab uno tantum recitátur Offícium, ante síngulas Lectiónes Matutíni atque ad Lectiónem brevem Primæ et Completórii, dícitur: Jube, Dómine, benedícere; et subjúngitur congruens Benedíctio. Ab Epíscopo autem, últimam Matutíni Lectiónem cantatúro, item dícitur: Jube, Dómine, benedícere; et respondétur a Choro: Amen.}

\variantheading{Pro I Lectione}

\benedlabel{Bened.} Benedictióne perpétua benedícat nos Pater ætérnus. \R\ Amen.

\rubric{Deinde dicuntur in unoquoque Nocturno Lectiones, prouti Officium occurrens requirit, et in fine cujuslibet Lectionis additur:}

\V\ Tu autem, Dómine, miserére nobis. \R\ Deo grátias.

\rubric{Post quamlibet vero Lectiónem, quæ Hymnum Te Deum immediáte non præcédat, cóngruens dícitur Responsórium, et in fine últimi Responsórii cujúsque Nocturni addítur Versus: Glória Patri, et Fílio, et Spirítui Sancto, et Responsórium ab Asterísco *, et quidem a secúndo, si duo fúerint, repétitur.}

\variantheading{Pro ii Lectione}

\benedlabel{Bened.} Unigénitus Dei Fílius nos benedícere et adjuváre dignétur. \R\ Amen.

\variantheading{Pro iii Lectione}

\benedlabel{Bened.} Spíritus Sancti grátia illúminet sensus et corda nostra. \R\ Amen.

\nocturnoheading{In II Nocturno}

\rubric{Sub congruentibus item Antiphonis dicuntur tres Psalmi et Versus, sicut in I Nocturno.}

\rubric{Post Versum dicitur.}

Pater noster \rubric{secreto usque ad} \V\ Et ne nos indúcas in tentatiónem. \R\ Sed líbera nos a malo.

\absolutioheading{Absolutio}

\lettrine[lines=2]{I}{PSIUS} píetas et misericórdia nos ádjuvet, qui cum Patre et Spíritu Sancto vivit et regnat in sæcula sæculórum. \R\ Amen.

\V\ Jube, domne, benedícere.

\variantheading{Pro iv Lectione}

\benedlabel{Bened.} Deus Pater omnípotens sit nobis propítius et clemens. \R\ Amen.

\variantheading{Pro v Lectione}

\benedlabel{Bened.} Christus perpétuæ det nobis gáudia vitæ. \R\ Amen.

\variantheading{Pro vi Lectione}

\benedlabel{Bened.} Ignem sui amóris accéndat Deus in córdibus nostris. \R\ Amen.

\nocturnoheading{In III Nocturno}

\rubric{Sub congruentibus denique Antiphonis dicuntur tres Psalmi et Versus, sicut in I et II Nocturno.}

\rubric{Post Versum dicitur}

Pater noster \rubric{secreto usque ad} \V\ Et ne nos indúcas in tentatiónem. \R\ Sed líbera nos a malo.

\absolutioheading{Absolutio}

\lettrine[lines=2]{A}{} \textsc{vinculis} peccatórum nostrórum absólvat nos omnípotens et miséricors Dóminus. \R\ Amen.

\V\ Jube, domne, benedícere.

\variantheading{Pro vii Lectione}

\benedlabel{Bened.} Evangélica léctio sit nobis salus et protéctio. \R\ Amen.

\variantheading{Pro viii Lectione}

\rubric{In Festis Domini et in Dominicis:}

\benedlabel{Bened.} Divínum auxílium máneat semper nobíscum. \R\ Amen.

\rubric{In Festis beatæ Mariæ Virginis:}

\benedlabel{Bened.} Cujus festum cólimus, ipsa Virgo virgínum intercédat pro nobis ad Dóminum. \R\ Amen.

\rubric{In Festis Sanctorum:}

\benedlabel{Bened.} Cujus (\textit{vel} Quorum \textit{aut} Quarum) festum cólimus, ipse (\textit{vel} ipsa \textit{aut} ipsi \textit{vel} ipsæ) intercédat (\textit{vel} intercédant) pro nobis ad Dóminum. \R\ Amen.

\variantheading{Pro ix Lectione}

\benedlabel{Bened.} Ad societátem cívium supernórum perdúcat nos Rex Angelórum. \R\ Amen.

\rubric{Si autem legénda sit Léctio ix de Homília cum Evangélio Domínicæ, vel Fériæ, aut Vigíliæ:}

\benedlabel{Bened.} Per evangélica dicta deleántur nostra delícta. \R\ Amen.

\variantheading{In Officio trium Lectionum}

\nocturnoheading{In Nocturno}

\rubric{In Officiis trium Lectionum, post Hymnum dicuntur Antiphonæ convenientes, quæ initio Psalmi inchoantur tantum et usque ad Asteriscum * perducuntur, ac deinde in fine integræ pronuntiantur.}

\rubric{Sub eisdem vero Antiphonis dicuntur novem Psalmi Feriæ currentis, quibus subjungitur Versus in III Nocturno positus, omissis Versibus pro I et II Nocturno assignatis.}

\rubric{Post Versum dicitur}

Pater noster \rubric{secreto usque ad} \V\ Et ne nos indúcas in tentatiónem. \R\ Sed líbera nos a malo.

\absolutioheading{Absolutio}

\rubric{Feria II et V}

Exáudi, Dómine, \rubric{ut in I Nocturno.}

\rubric{Feria III et VI}

Ipsíus píetas, \rubric{ut in II Nocturno.}

\rubric{Feria IV et Sabbato}

A vínculis, \rubric{ut in III Nocturno Officii novem Lectionum.}

\rubric{Deinde leguntur Lectiones cum Responsoriis, prouti Officium occurrens requirit; et ante eas dicuntur sequentes}

\variantheading{Benedictiones}

\rubric{In Feriis, quando legitur Homilia cum Evangelio, erit:}

\benedlabel{1 Bened.} Evangélica léctio sit nobis salus et protéctio. \R\ Amen.

\benedlabel{2 Bened.} Divínum auxílium máneat semper nobíscum. \R\ Amen.

\benedlabel{3 Bened.} Ad societátem cívium supernórum perdúcat nos Rex Angelórum. \R\ Amen.

\rubric{In Feriis, quando non legitur Homilia cum Evangelio, Feria II et V Benedictiones ut in I Nocturno, Feria III et VI ut in II Nocturno Officii novem Lectionum; Feria autem IV et Sabbato erit:}

\benedlabel{1 Bened.} Ille nos benedícat, qui sine fine vivit et regnat. \R\ Amen.

\benedlabel{2 Bened.} Divínum auxílium máneat semper nobíscum. \R\ Amen.

\benedlabel{3 Bened.} Ad societátem cívium supernórum perdúcat nos Rex Angelórum. \R\ Amen.

\rubric{In Festis Sanctorum erit:}

\benedlabel{1 Bened.} Ille nos benedícat, qui sine fine vivit et regnat. \R\ Amen.

\benedlabel{2 Bened.} Cujus (\textit{vel} Quorum \textit{aut} Quarum) festum cólimus, ipse (\textit{vel} ipsa \textit{aut} ipsi \textit{vel} ipsæ) intercédat (\textit{vel} intercédant) pro nobis ad Dóminum. \R\ Amen.

\benedlabel{3 Bened.} Ad societátem cívium supernórum perdúcat nos Rex Angelórum. \R\ Amen.

\rubric{In Officio de sancta María in Sábbato Absolútio et Benedictiónes, quæ ibídem pónuntur.}

\rubric{Post últimam Lectiónem, in ómnibus Domínicis Témporis Paschális, in Festis cujúsvis ritus, per omnes Octávas, et in Offício sanctæ Maríæ in Sábbato, dícitur Hymnus Ambrosiánus. A Domínica autem I Quadragésimæ usque ad Sábbatum sanctum inclusíve, non dícitur nisi in Festis; a Paschate vero usque ad Pentecósten inclusíve, dícitur étiam in feriáli Offício, excépta Féria II Rogatiónum. Quando vero Hymnus prædíctus omíttitur, ejus loco dícitur IX aut III Responsórium.}

\hymnusheading{Hymnus Ambrosianus}

\lettrine[lines=2]{T}{E} Deum laudámus: * te Dóminum confitémur.

Te ætérnum Patrem * omnis terra venerátur.

Tibi omnes Angeli, * tibi Cæli, et univérsæ Potestátes:

Tibi Chérubim et Séraphim * incessábili voce proclámant:

Sanctus, * Sanctus, * Sanctus * Dóminus, Deus Sábaoth.

Pleni sunt cæli et terra * majestátis glóriæ tuæ.

Te gloriósus * Apostolórum chorus,

Te Prophetárum * laudábilis númerus,

Te Mártyrum candidátus * laudat exércitus.

Te per orbem terrárum * sancta confitétur Ecclésia,

Patrem * imménsæ majestátis;

Venerándum tuum verum * et únicum Fílium;

Sanctum quoque * Paráclitum Spíritum.

Tu Rex glóriæ, * Christe.

Tu Patris * sempitérnus es Fílius.

Tu, ad liberándum susceptúrus hóminem, * non horruísti Vírginis úterum.

Tu, devícto mortis acúleo, * aperuísti credéntibus regna cælórum.

Tu ad déxteram Dei sedes, * in glória Patris.

Judex créderis * esse ventúrus.

\rubric{¶ Sequens versus dicitur flexis genibus.}

Te ergo quæsumus, tuis fámulis súbveni, * quos pretióso sánguine redemísti.

Ætérna fac cum Sanctis tuis * in glória numerári.

Salvum fac pópulum tuum, Dómine, * et bénedic hereditáti tuæ.

Et rege eos, * et extólle illos usque in ætérnum.

Per síngulos dies * benedícimus te;

Et laudámus nomen tuum in séculum, * et in séculum séculi.

Dignáre, Dómine, die isto * sine peccáto nos custodíre.

Miserére nostri, Dómine, * miserére nostri.

Fiat misericórdia tua, Dómine, super nos, * quemádmodum sperávimus in te.

In te, Dómine, sperávi: * non confúndar in ætérnum.

\rubric{Dicto Te Deum, aut ultimo Responsorio, statim incipiuntur Laudes a Versu Deus, in adjutorio, ut infra.}

\rubric{In privata vero recitatione Matutinum a Laudibus separari potest; quo in casu, post Hymnum Te Deum, vel post ultimum Responsorium, dicitur:}

\V\ Dóminus vobíscum. \R\ Et cum spíritu tuo.

\rubric{¶ Hic Versus Dóminus vobíscum non dícitur ab eo, qui non est saltem in órdine Diaconátus; sed ejus loco substitúitur, étiam in áliis Horis, quóties in antecedéntibus Précibus dictus non fúerit:}

\V\ Dómine, exáudi oratiónem meam. \R\ Et clamor meus ad te véniat.

Orémus.

\rubric{Et dicitur Oratio ut ad Laudes.}

\rubric{Post Orationem additur:}

\V\ Dóminus vobíscum. \R\ Et cum spíritu tuo.

\V\ Benedicámus Dómino. \R\ Deo grátias.

\V\ Fidélium ánimæ per misericórdiam Dei requiéscant in pace. \R\ Amen.

\rubric{Deinde dicitur Pater noster totum secreto.}

\rubric{In publica vero recitatione Matutinum in Choro a Laudibus separari non potest.}

\hourheading{Ad Laudes}

\rubric{Si Laudes extra Chorum separéntur a Matutíno, ante eas dícitur secréto Pater noster et Ave María; secus absolúte incipiúntur, ut séquitur:}

\V\ Deus in adjutórium.

\rubric{Deinde, sub congruéntibus Antiphónis, dicúntur Psalmi et Cánticum, prouti Offícium occúrrens requírit.}

\rubric{Repetita vero post últimum Psalmum Antíphona, dicúntur Capítulum, Hymnus et Versus, prouti Offícium occúrrens exígit.}

\variantheading{Tempore Quadragesimæ}

\rubric{In Officio dominicali, Capitulum ut in Proprio de Tempore.}

\rubric{In Officio autem feriali dicitur quotidie sequens}

\capitulumlabel{Capitulum}{Is. 58, 1}

\lettrine[lines=2]{C}{LAMA}, ne cesses, quasi tuba exálta vocem tuam, et annúntia pópulo meo scélera eórum, et dómui Jacob peccáta eórum. \R\ Deo grátias.

\rubric{¶ Sic respondetur in fine omnium Capitulorum.}

\rubric{In Officio tamen dominicali et feriali dicitur sequens Hymnus et Versus:}

\hymnusheading{Hymnus}

\lettrine[lines=2]{O}{} Sol salútis, íntimis,\\
Jesu, refúlge méntibus,\\
Dum, nocte pulsa, grátior\\
Orbi dies renáscitur.

Dans tempus acceptábile,\\
Da lacrimárum rívulis\\
Laváre cordis víctimam,\\
Quam læta adúrat cáritas.

Quo fonte manávit nefas,\\
Fluent perénnes lácrimæ,\\
Si virga pæniténtiæ\\
Cordis rigórem cónterat.

Dies venit, dies tua,\\
In qua reflórent ómnia:\\
Lætémur et nos, in viam\\
Tua redúcti déxtera.

\rubric{¶ Sequens Conclusio numquam mutatur.}

Te prona mundi máchina,\\
Clemens, adóret, Trínitas,\\
Et nos novi per grátiam\\
Novum canámus cánticum. Amen.

\V\ Angelis suis Deus mandávit de te. \R\ Ut custódiant te in ómnibus viis tuis.

\variantheading{Tempore Passionis}

\rubric{In Officio dominicali, Capitulum ut in Proprio de Tempore.}

\rubric{In Officio autem feriali, usque ad Feriam IV Majoris Hebdomadæ inclusive, dicitur quotidie sequens}

\capitulumlabel{Capitulum}{Jer. 11, 19}

\lettrine[lines=2]{V}{ENITE}, mittámus lignum in panem ejus, et eradámus eum de terra vivéntium, et nomen ejus non memorétur ámplius.

\rubric{In Officio tamen dominicali et feriali dicitur sequens Hymnus et Versus:}

\hymnusheading{Hymnus}

\lettrine[lines=2]{L}{USTRA} sex qui jam perégit,\\
Tempus implens córporis,\\
Sponte líbera Redémptor\\
Passióni déditus,\\
Agnus in Crucis levátur\\
Immolándus stípite.

Felle potus ecce languet:\\
Spina, clavi, láncea\\
Mite corpus perforárunt:\\
Unda manat, et cruor:\\
Terra, pontus, astra, mundus,\\
Quo lavántur flúmine!

Crux fidélis, inter omnes\\
Arbor una nóbilis:\\
Silva talem nulla profert\\
Fronde, flore, gérmine:\\
Dulce ferrum, dulce lignum,\\
Dulce pondus sústinent.

Flecte ramos, arbor alta,\\
Tensa laxa víscera,\\
Et rigor lentéscat ille,\\
Quem dedit natívitas;\\
Et supérni membra Regis\\
Tende miti stípite.

Sola digna tu fuísti\\
Ferre mundi víctimam,\\
Atque portum præparáre\\
Arca mundo náufrago,\\
Quam sacer cruor perúnxit,\\
Fusus Agni córpore.

Sempitérna sit beátæ\\
Trinitáti glória,\\
Æqua Patri, Filióque;\\
Par decus Paráclito:\\
Uníus Triníque nomen\\
Laudet univérsitas. Amen.

\V\ Eripe me de inimícis meis, Deus meus. \R\ Et ab insurgéntibus in me líbera me.

\variantheading{Tempore Paschali}

\rubric{In Officio dominicali, a Dominica in Albis usque ad Dominicam V post Pascha inclusive, Capitulum ut in Proprio de Tempore.}

\rubric{In Officio autem feriali, a Feria II post Dominicam in Albis usque ad Vigiliam Ascensionis inclusive, dicitur quotidie sequens}

\capitulumlabel{Capitulum}{Rom. 6, 9-10}

\lettrine[lines=2]{C}{HRISTUS} resúrgens ex mórtuis jam non móritur, mors illi ultra non dominábitur. Quod enim mórtuus est peccáto, mórtuus est semel: quod autem vivit, vivit Deo.

\rubric{In Officio tamen dominicali et feriali dicitur sequens Hymnus et Versus:}

\hymnusheading{Hymnus}

\lettrine[lines=2]{A}{URORA} cælum púrpurat,\\
Æther resúltat láudibus,\\
Mundus triúmphans júbilat,\\
Horrens avérnus ínfremit:

Rex ille dum fortíssimus\\
De mortis inférno specu\\
Patrum senátum líberum\\
Edúcit ad vitæ jubar.

Cujus sepúlcrum plúrimo\\
Custóde signábat lapis,\\
Victor triúmphat, et suo\\
Mortem sepúlcro fúnerat.

Sat fúneri, sat lácrimis,\\
Sat est datum dolóribus:\\
Surréxit exstínctor necis,\\
Clamat corúscans Angelus.

Ut sis perénne méntibus\\
Paschále, Jesu, gáudium,\\
A morte dira críminum\\
Vitæ renátos líbera.

Deo Patri sit glória,\\
Et Fílio, qui a mórtuis\\
Surréxit, ac Paráclito,\\
In sempitérna sæcula. Amen.

\V\ In resurrectióne tua, Christe, allelúja. \R\ Cæli et terra læténtur, allelúja.

\rubric{Recitato Versu, dicitur cum Antiphona convenienti sequens}

\canticumheading{Canticum Zachariæ}{Luc. 1, 68-69}

\lettrine[lines=2]{B}{ENEDICTUS} \textdagger\ Dóminus, Deus Israël, * quia visitávit, et fecit redemptiónem plebis suæ:

Et eréxit cornu salútis nobis: * in domo David, púeri sui.

Sicut locútus est per os sanctórum, * qui a século sunt, Prophetárum ejus:

Salútem ex inimícis nostris, * et de manu ómnium, qui odérunt nos:

Ad faciéndam misericórdiam cum pátribus nostris: * et memorári testaménti sui sancti.

Jusjurándum, quod jurávit ad Abraham, patrem nostrum, * datúrum se nobis:

Ut sine timóre, de manu inimicórum nostrórum liberáti, * serviámus illi.

In sanctitáte, et justítia coram ipso, * ómnibus diébus nostris.

Et tu, puer, Prophéta Altíssimi vocáberis: * præíbis enim ante fáciem Dómini paráre vias ejus:

Ad dandam sciéntiam salútis plebi ejus: * in remissiónem peccatórum eórum:

Per víscera misericórdiæ Dei nostri: * in quibus visitávit nos, óriens ex alto:

Illumináre his, qui in ténebris, et in umbra mortis sedent: * ad dirigéndos pedes nostros in viam pacis.

Glória Patri, et Fílio, * et Spirítui Sancto.

Sicut erat in princípio, et nunc, et semper, * et in sécula seculórum. Amen.

\rubric{¶ Hic Versus Glória Patri dicitur in fine omnium Canticorum, nisi aliter adnotetur.}

\rubric{Expleto Cantico ac repetita Antiphona, si Preces non fuerint recitandæ, statim dicitur Versus Dóminus vobíscum cum Oratione et reliquis omnibus usque ad finem Horæ, ut infra.}

\rubric{In Feriis vero Quadragesimæ, et Passionis usque ad Feriam IV Majoris Hebdomadæ inclusive, si fiat Officium de Feria currenti, et quamvis in eo peragenda sit Commemoratio Duplicis aut Semiduplicis ad instar Simplicis redacti, post repetitam ad Benedictus Antiphonam dicuntur flexis genibus sequentes Preces, quæ aliis temporibus omittuntur:}

KYRIE, eléison. Christe, eléison. Kýrie, eléison.

\rubric{Pater noster, quod dicitur a solo Hebdomadario totum clara voce, usque ad} \V\ Et ne nos indúcas in tentatiónem. \R\ Sed líbera nos a malo.

\V\ Ego dixi: Dómine, miserére mei. \R\ Sana ánimam meam, quia peccávi tibi. \V\ Convértere, Dómine, úsquequo? \R\ Et deprecábilis esto super servos tuos. \V\ Fiat misericórdia tua, Dómine, super nos. \R\ Quemádmodum sperávimus in te. \V\ Sacerdótes tui induántur justítiam. \R\ Et sancti tui exsúltent.

\V\ Orémus pro beatíssimo Papa nostro N. \R\ Dóminus consérvet eum, et vivíficet eum, et beátum fáciat eum in terra, et non tradat eum in ánimam inimicórum ejus.

\V\ Orémus et pro Antístite nostro N. \R\ Stet et pascat in fortitúdine tua, Dómine, in sublimitáte nóminis tui.

\rubric{¶ Romæ præcedens Versus cum suo Responsorio omittitur; alibi vero, ad litteram N., ab omnibus nomen diœcesani Episcopi exprimatur. Vacante Apostolica vel Episcopali Sede, alteruter vel uterque respectivus Versus cum suo Responsorio prætéritur.}

\V\ Dómine, salvum fac regem. \R\ Et exáudi nos in die, qua invocavérimus te. \V\ Salvum fac pópulum tuum, Dómine, et bénedic hereditáti tuæ. \R\ Et rege eos, et extólle illos usque in ætérnum. \V\ Meménto Congregatiónis tuæ. \R\ Quam possedísti ab inítio. \V\ Fiat pax in virtúte tua. \R\ Et abundántia in túrribus tuis. \V\ Orémus pro benefactóribus nostris. \R\ Retribúere dignáre, Dómine, ómnibus, nobis bona faciéntibus propter nomen tuum, vitam ætérnam. Amen. \V\ Orémus pro fidélibus defúnctis. \R\ Réquiem ætérnam dona eis, Dómine, et lux perpétua lúceat eis. \V\ Requiéscant in pace. \R\ Amen. \V\ Pro frátribus nostris abséntibus. \R\ Salvos fac servos tuos, Deus meus, sperántes in te. \V\ Pro afflíctis et captívis. \R\ Líbera eos, Deus Israël, ex ómnibus tribulatiónibus suis. \V\ Mitte eis, Dómine, auxílium de sancto. \R\ Et de Sion tuére eos. \V\ Dómine, Deus virtútum, convérte nos. \R\ Et osténde fáciem tuam, et salvi érimus. \V\ Exsúrge, Christe, ádjuva nos. \R\ Et líbera nos propter nomen tuum. \V\ Dómine, exáudi oratiónem meam. \R\ Et clamor meus ad te véniat.

\rubric{Absolutis Precibus, vel, si Preces locum non habeant, repetita Antiphona post Benedictus, statim dicitur:}

\V\ Dóminus vobíscum. \R\ Et cum spíritu tuo.

Orémus.

\rubric{Et dicitur Oratio conveniens. Postea fiunt Commemorationes, si occurrant.}

\rubric{Deinde, extra Tempus Paschale, in Officio semiduplici aut simplici tam de Tempore quam de Sanctis, præter quam a Dominica Passionis usque ad Feriam IV Majoris Hebdomadæ inclusive, ac præter omnes dies, in quibus habita sit Commemoratio Duplicis, fit sequens}

\suffragiumheading{Suffragium de omnibus Sanctis}

\antlabel{Ant.} Beáta Dei Génitrix Virgo María, Sanctíque omnes intercédant pro nobis ad Dóminum.

\V\ Mirificávit Dóminus Sanctos suos. \R\ Et exaudívit eos clamántes ad se.

Orémus. \oratiolabel{Oratio}

\lettrine[lines=2]{A}{} \textsc{cunctis} nos, quæsumus, Dómine, mentis et córporis defénde perículis: et, intercedénte beáta et gloriósa semper Vírgine Dei Genitríce María, cum beáto Joseph, beátis Apóstolis tuis Petro et Paulo, atque beáto N. et ómnibus Sanctis, salútem nobis tríbue benígnus et pacem; ut, destrúctis adversitátibus et erróribus univérsis, Ecclésia tua secúra tibi sérviat libertáte. Per eúmdem Dóminum nostrum. \R\ Amen.

\rubric{¶ In hac Oratione, ad litteram N., exprimitur nomen Titularis propriæ Ecclesiæ, dummodo Titulus non sit Persona divina vel Mysterium Domini, aut de ipso ratione Vigiliæ vel Festi factum non fuerit Officium vel Commemoratio, aut illius nomen non sit in eadem Suffragii Oratione expressum; ac nomina sanctorum Angelorum et sancti Joannis Baptistæ, si Titulares fuerint, præponuntur nomini sancti Joseph. In hisce autem omnibus casibus omittuntur verba atque beáto N.}

\rubric{Tempore autem Paschali, a Feria II post Dominicam in Albis usque ad Vigiliam Ascensionis inclusive, in Officio Semiduplici aut Simplici tam de Tempore quam de Sanctis, nisi habita sit Commemoratio Duplicis, vel factum fuerit Officium aut Commemoratio cujusvis Octavæ, loco Suffragii, fit sequens}

\suffragiumheading{Commemoratio de Cruce}

\antlabel{Ant.} Crucifíxus surréxit a mórtuis, et redémit nos, allelúja, allelúja.

\V\ Dícite in natiónibus, allelúja. \R\ Quia Dóminus regnávit a ligno, allelúja.

Orémus. \oratiolabel{Oratio}

\lettrine[lines=2]{D}{EUS}, qui pro nobis Fílium tuum Crucis patíbulum subíre voluísti, ut inimíci a nobis expélleres potestátem: concéde nobis, fámulis tuis; ut resurrectiónis grátiam consequámur. Per eúmdem Dóminum. \R\ Amen.

\rubric{Post ultimam Orationem additur:}

\V\ Dóminus vobíscum. \R\ Et cum spíritu tuo. \V\ Benedicámus Dómino. \R\ Deo grátias. \V\ Fidélium ánimæ per misericórdiam Dei requiéscant in pace. \R\ Amen.

\rubric{Deinde, in publica Officii recitatione, si discedendum sit a Choro, itemque in recitatione privata, si tunc terminetur Officium; alioquin, in fine ultimæ Horæ dicitur Pater noster, totum secreto; eoque recitato, subjungitur:}

\V\ Dóminus det nobis suam pacem. \R\ Et vitam ætérnam. Amen.

\rubric{Et immediate dicitur, cum suis Versu et Oratione, una ex finalibus beatæ Mariæ Virginis Antiphonis, quæ inferius post Completorium pro diversitate Temporis assignantur.}

\rubric{Postea concluditur:}

\V\ Divínum auxílium máneat semper nobíscum. \R\ Amen.

\hourheading{Ad Primam}

Pater, Ave et Credo.

\V\ Deus, in adjutórium.

\hymnusheading{Hymnus}

\lettrine[lines=2]{J}{AM} lucis orto sídere,\\
Deum precémur súpplices,\\
Ut in diúrnis áctibus\\
Nos servet a nocéntibus.

Linguam refrénans témperet,\\
Ne litis horror ínsonet:\\
Visum fovéndo cóntegat,\\
Ne vanitátes háuriat.

Sint pura cordis íntima,\\
Absístat et vecórdia:\\
Carnis terat supérbiam\\
Potus cibíque párcitas.

Ut, cum dies abscésserit\\
Noctémque sors redúxerit,\\
Mundi per abstinéntiam\\
Ipsi canámus glóriam.

Deo Patri sit glória,\\
Ejúsque soli Fílio,\\
Cum Spíritu Paráclito,\\
Nunc et per omne sæculum. Amen.

\rubric{¶ Conclusio communis, in Hymno præcedenti, et in aliis idem metrum habentibus, semper omittitur, quando specialis in omnibus Horis adhibenda præscribitur; et, si plures Conclusiones propriæ occurrant, sumitur Conclusio Officii currentis, aut secus Officii ipsa die et primo quidem loco inter cetera propriam Conclusionem habentia commemorandi, aut demum de occurrenti Octava communi vel de Tempore.}

\rubric{Expleto Hymno, convéniens dicitur Antiphona usque ad Asteriscum *, prouti occurrens Officium requirit.}

\variantheading{Tempore Quadragesimæ}

\rubric{In Officio dominicali, Antiphona ut in Proprio de Tempore.}

\rubric{In Officio autem feriali, dicitur quotidie sequens}

\antlabel{Ant.} Vivo ego, * dicit Dóminus: nolo mortem peccatóris, sed ut magis convertátur, et vivat.

\variantheading{Tempore Passionis}

\rubric{In Officio dominicali, Antiphona ut in Proprio de Tempore.}

\rubric{In Officio feriali: In omnibus Feriis infra hebdomadam Passionis dicitur sequens}

\antlabel{Ant.} Líbera me, Dómine, * et pone me juxta te: et cujúsvis manus pugnet contra me.

\rubric{In Feria vero II, III et IV Majóris Hebdómadæ Antiphona prima ex Laudibus Fériæ curréntis ut in Proprio de Témpore.}

\variantheading{Tempore Paschali}

\rubric{In Officio Dominicali et feriali, a Dominica in Albis usque ad Vigiliam Ascensionis, in Officio autem de Sanctis usque ad Feriam VI post Octavam Ascensionis inclusive, Antiphona} Allelúja, \rubric{etiam ad Tertiam, Sextam, Nonam et Completorium, ut in Psalterio pro singulis Feriis apponitur.}

\rubric{Post enuntiatam Antiphonam dicuntur tres Psalmi Officio diei congruentes.}

\rubric{¶ Quando autem adhibetur posterius Laudum schema, duo Psalmi, in priori Laudum dominicalium schemate omissi, resumuntur ad Primam Dominicæ ante alios Psalmos; et unus Psalmus, in priori pariter Laudum ferialium schemate omissus, item ad Primam Feriæ currentis post alios Psalmos resumitur.}

\rubric{Repetita post ultimum Psalmum integre Antiphona, conveniens dicitur Capitulum; nimirum:}

\rubric{In omnibus Dominicis, in Vigilia privilegiata Pentecostes, in Officio cujuslibet Festi vel Octavæ, ac sanctæ Mariæ in Sabbato, necnon in Feriis Temporis Paschalis:}

\capitulumlabel{Capitulum}{1 Tim. 1, 17}

\lettrine[lines=2]{R}{EGI} sæculórum immortáli et invisíbili, soli Deo honor et glória in sæcula sæculórum. Amen.

\rubric{In omnibus autem Feriis extra Tempus Paschale:}

\capitulumlabel{Capitulum}{Zach. 8, 19}

\lettrine[lines=2]{P}{ACEM} et veritátem dilígite, ait Dóminus omnípotens.

\rubric{Dicto Capitulo, subjungitur Responsorium breve, ut sequitur. Versus tamen Qui sedes ad déxteram Patris semper omittitur, quando specialis Versus adhibendus præscribitur, Officii nempe currentis, aut secus Officii ad Laudes et primo quidem loco inter cetera specialem Versum habentia commemorati, aut demum de occurrenti Octava communi vel de Tempore proprius.}

\R\ br. Christe, Fili Dei vivi, * Miserére nobis. \rubric{Et repetitur:} Christe, Fili Dei vivi, miserére nobis. \V\ Qui sedes ad déxteram Patris. \R\ Miserére nobis. \V\ Glória Patri, et Fílio, et Spirítui Sancto. \R\ Christe, Fili Dei vivi, miserére nobis. \V\ Exsúrge, Christe, ádjuva nos. \R\ Et líbera nos propter nomen tuum.

\rubric{¶ Ita dicuntur semper Responsoria brevia, etiam ad Tertiam, Sextam, Nonam et Completorium.}

\rubric{Tempore Passionis, in Officio dominicali et feriali, omittitur in Responsorio brevi Glória Patri, et hujus loco repetitur Responsorium, ut sequitur:}

Christe. \V\ \textit{Qui sedes ad déxteram Patris.} Miserére nobis. Christe, Fili Dei vivi. Miserére nobis.

\V\ Exsúrge, Christe, ádjuva nos. \R\ Et líbera nos propter nomen tuum.

\rubric{Tempore Paschali dicitur sequens}

\R\ br. Christe, Fili Dei vivi, miserére nobis, * Allelúja, allelúja. Christe. \V\ \textit{Qui surrexísti a mórtuis.} Allelúja, allelúja. Glória Patri. Christe.

\V\ Exsúrge, Christe, ádjuva nos, allelúja. \R\ Et líbera nos propter nomen tuum, allelúja.

\rubric{Et sic dicitur Responsorium breve a Dominica in Albis usque ad Sabbatum infra Octavam Pentecostes inclusive, etiam in Festis; excepto Versu} Qui surrexísti a mórtuis, \rubric{si alter Versus in Proprio de Tempore vel de Sanctis assignetur.}

\rubric{Expleto Responsorio brevi, si Preces non fuerint recitandæ, statim dicitur Versus} Dóminus vobíscum \rubric{cum subsequenti Oratione et reliquis omnibus usque ad finem Horæ ut infra.}

\rubric{In quolibet Officio semiduplici, in Officio sanctæ Mariæ in Sabbato, in Festis simplicibus, atque in Feriis Temporis Paschalis post Responsorium breve dicuntur sequentes Preces dominicales; quæ tamen omittuntur in Vigília Pentecóstes et ómnibus ínsuper diébus, in quibus ad Laudes habita sit Commemoratio Duplicis, vel factum fuerit Officium aut Commemoratio cujusvis Octavæ vel Feriæ VI post Octavam Ascensionis:}

KYRIE, eléison. Christe, eléison. Kýrie, eléison.

Pater noster \rubric{secreto usque ad} \V\ Et ne nos indúcas in tentatiónem. \R\ Sed líbera nos a malo.

Credo in Deum \rubric{secreto usque ad} \V\ Carnis resurrectiónem. \R\ Vitam ætérnam. Amen.

\V\ Et ego ad te, Dómine, clamávi. \R\ Et mane orátio mea prævéniet te. \V\ Repleátur os meum laude. \R\ Ut cantem glóriam tuam, tota die magnitúdinem tuam. \V\ Dómine, avérte fáciem tuam a peccátis meis. \R\ Et omnes iniquitátes meas dele. \V\ Cor mundum crea in me, Deus. \R\ Et spíritum rectum ínnova in viscéribus meis. \V\ Ne projícias me a fácie tua. \R\ Et spíritum sanctum tuum ne áuferas a me. \V\ Redde mihi lætítiam salutáris tui. \R\ Et spíritu principáli confírma me. \V\ Adjutórium nostrum in nómine Dómini. \R\ Qui fecit cælum et terram.

\rubric{Deinde Hebdomadarius facit Confessionem cum Absolutione, eamque Chorus repetit, ut infra post feriales Preces.}

\rubric{In Feriis vero Quadragesimæ, et Passionis usque ad Feriam IV Majoris Hebdomadæ inclusive, si ad Laudes recitatæ fuerint Preces, post Responsorium breve dicuntur flexis genibus sequentes Preces feriales:}

KYRIE, eléison. Christe, eléison. Kýrie, eléison.

Pater noster \rubric{secreto usque ad} \V\ Et ne nos indúcas in tentatiónem. \R\ Sed líbera nos a malo.

Credo in Deum \rubric{secreto usque ad} \V\ Carnis resurrectiónem. \R\ Vitam ætérnam. Amen.

\V\ Et ego ad te, Dómine, clamávi. \R\ Et mane orátio mea prævéniet te. \V\ Repleátur os meum laude. \R\ Ut cantem glóriam tuam, tota die magnitúdinem tuam. \V\ Dómine, avérte fáciem tuam a peccátis meis. \R\ Et omnes iniquitátes meas dele. \V\ Cor mundum crea in me, Deus. \R\ Et spíritum rectum ínnova in viscéribus meis. \V\ Ne projícias me a fácie tua. \R\ Et spíritum sanctum tuum ne áuferas a me. \V\ Redde mihi lætítiam salutáris tui. \R\ Et spíritu principáli confírma me. \V\ Eripe me, Dómine, ab hómine malo. \R\ A viro iníquo éripe me. \V\ Eripe me de inimícis meis, Deus meus. \R\ Et ab insurgéntibus in me líbera me. \V\ Eripe me de operántibus iniquitátem. \R\ Et de viris sánguinum salva me. \V\ Sic psalmum dicam nómini tuo in sæculum sæculi. \R\ Ut reddam vota mea de die in diem. \V\ Exáudi nos, Deus, salutáris noster. \R\ Spes ómnium fínium terræ, et in mari longe. \V\ Deus, in adjutórium meum inténde. \R\ Dómine, ad adjuvándum me festína. \V\ Sanctus Deus, Sanctus fortis, Sanctus immortális. \R\ Miserére nobis. \V\ Bénedic, ánima mea, Dómino. \R\ Et ómnia, quæ intra me sunt, nómini sancto ejus. \V\ Bénedic, ánima mea, Dómino. \R\ Et noli oblivísci omnes retributiónes ejus. \V\ Qui propitiátur ómnibus iniquitátibus tuis. \R\ Qui sanat omnes infirmitátes tuas. \V\ Qui rédimit de intéritu vitam tuam. \R\ Qui corónat te in misericórdia et miseratiónibus. \V\ Qui replet in bonis desidérium tuum. \R\ Renovábitur ut áquilæ juvéntus tua. \V\ Adjutórium nostrum in nómine Dómini. \R\ Qui fecit cælum et terram.

\rubric{Deinde Hebdomadarius facit Confessionem:}

\lettrine[lines=2]{C}{ONFITEOR} Deo omnipoténti, beátæ Maríæ semper Vírgini, beáto Michaéli Archángelo, beáto Joánni Baptístæ, sanctis Apóstolis Petro et Paulo, ómnibus Sanctis, et vobis, fratres, quia peccávi nimis cogitatióne, verbo et ópere: mea culpa, mea culpa, mea máxima culpa. Ideo precor beátam Maríam semper Vírginem, beátum Michaélem Archángelum, beátum Joánnem Baptístam, sanctos Apóstolos Petrum et Paulum, omnes Sanctos, et vos, fratres, oráre pro me ad Dóminum, Deum nostrum.

\rubric{Chorus respondet:}

\lettrine[lines=2]{M}{ISEREATUR} tui omnípotens Deus, et, dimíssis peccátis tuis, perdúcat te ad vitam ætérnam. \R\ Amen.

\rubric{Deinde repetit Confessionem, et ubi dicitur: vobis, fratres, et vos, fratres, dicat: tibi, pater, et te, pater.}

\rubric{Facta Confessione a Choro, Hebdomadarius dicit:}

\lettrine[lines=2]{M}{ISEREATUR} vestri omnípotens Deus, et, dimíssis peccátis vestris, perdúcat vos ad vitam ætérnam. \R\ Amen.

\lettrine[lines=2]{I}{NDULGENTIAM}, absolutiónem, et remissiónem peccatórum nostrórum tríbuat nobis omnípotens et miséricors Dóminus. \R\ Amen.

\rubric{Extra Chorum, si unus vel duo tantum recitent Officium, et in Choro Monialium, semel tantum ac simul ab omnibus ita fit Confessio:}

\lettrine[lines=2]{C}{ONFITEOR} Deo omnipoténti, beátæ Maríæ semper Vírgini, beáto Michaéli Archángelo, beáto Joánni Baptístæ, sanctis Apóstolis Petro et Paulo, et ómnibus Sanctis, quia peccávi nimis cogitatióne, verbo et ópere: mea culpa, mea culpa, mea máxima culpa. Ideo precor beátam Maríam semper Vírginem, beátum Michaélem Archángelum, beátum Joánnem Baptístam, sanctos Apóstolos Petrum et Paulum, et omnes Sanctos, oráre pro me ad Dóminum, Deum nostrum.

\rubric{Deinde dicitur:}

\lettrine[lines=2]{M}{ISEREATUR} nostri omnípotens Deus, et, dimíssis peccátis nostris, perdúcat nos ad vitam ætérnam. Amen.

\lettrine[lines=2]{I}{NDULGENTIAM}, absolutiónem, et remissiónem peccatórum nostrórum tríbuat nobis omnípotens et miséricors Dóminus. Amen.

\rubric{Facta Absolutione, subjungitur:}

\V\ Dignáre, Dómine, die isto. \R\ Sine peccáto nos custodíre. \V\ Miserére nostri, Dómine. \R\ Miserére nostri. \V\ Fiat misericórdia tua, Dómine, super nos. \R\ Quemádmodum sperávimus in te. \V\ Dómine, exáudi oratiónem meam. \R\ Et clamor meus ad te véniat.

\rubric{Absolutis Precibus, vel, si Preces locum non habeant, post Responsorium breve statim dicitur:}

\V\ Dóminus vobíscum. \R\ Et cum spíritu tuo.

Orémus. \oratiolabel{Oratio}

\lettrine[lines=2]{D}{OMINE}, Deus omnípotens, qui ad princípium hujus diéi nos perveníre fecísti: tua nos hódie salva virtúte; ut in hac die ad nullum declinémus peccátum, sed semper ad tuam justítiam faciéndam nostra procédant elóquia, dirigántur cogitatiónes et ópera. Per Dóminum nostrum. \R\ Amen.

\V\ Dóminus vobíscum. \R\ Et cum spíritu tuo. \V\ Benedicámus Dómino. \R\ Deo grátias.

\rubric{Deinde in Choro legitur Martyrologium, quod laudabiliter fit etiam extra Chorum. Postea Hebdomadarius dicit:}

\V\ Pretiósa in conspéctu Dómini. \R\ Mors Sanctórum ejus.

\rubric{Et absolute, sine} Orémus:

\oratiolabel{Oratio}

\lettrine[lines=2]{S}{ANCTA} María, et omnes Sancti intercédant pro nobis ad Dóminum, ut nos mereámur ab eo adjuvári et salvári, qui vivit et regnat in sæcula sæculórum. \R\ Amen.

\V\ Deus, in adjutórium meum inténde. \R\ Dómine, ad adjuvándum me festína. \rubric{Et dicitur ter; ultimo additur:} Glória Patri, et Fílio, et Spirítui Sancto. Sicut erat in princípio, et nunc, et semper, et in sæcula sæculórum. Amen.

Kýrie, eléison. Christe, eléison. Kýrie, eléison.

Pater noster \rubric{secreto usque ad} \V\ Et ne nos indúcas in tentatiónem. \R\ Sed líbera nos a malo.

\V\ Réspice in servos tuos, Dómine, et in ópera tua, et dírige fílios eórum. \R\ Et sit splendor Dómini, Dei nostri, super nos, et ópera mánuum nostrárum dírige super nos, et opus mánuum nostrárum dírige. \V\ Glória Patri. \R\ Sicut erat.

Orémus. \oratiolabel{Oratio}

\lettrine[lines=2]{D}{IRIGERE} et sanctificáre, régere et gubernáre dignáre, Dómine Deus, Rex cæli et terræ, hódie corda et córpora nostra, sensus, sermónes et actus nostros in lege tua, et in opéribus mandatórum tuórum: ut hic, et in ætérnum, te auxiliánte, salvi et líberi esse mereámur, Salvátor mundi: Qui vivis et regnas in sæcula sæculórum. \R\ Amen.

\rubric{Deinde dicitur Lectio brevis, ad absolutionem Capituli, prouti Officium occurrens requirit.}

\V\ Jube, domne, benedícere. \benedlabel{Bened.} Dies et actus nostros in sua pace dispónat Dóminus omnípotens. \R\ Amen.

\variantheading{In Officio dominicali et feriali}

\rubric{A Dominica I Quadragesimæ usque ad Sabbatum ante Dominicam Passionis inclusive:}

\lectioheading{Lectio brevis}{Is. 55, 6}

\lettrine[lines=2]{Q}{UÆRITE} Dóminum, dum inveníri potest: invocáte eum, dum prope est. Tu autem, Dómine, miserére nobis. \R\ Deo grátias.

\rubric{¶ Ita semper concluduntur Lectiones breves ad Primam, et ita in earum fine respondetur.}

\rubric{A Dominica Passionis usque ad Feriam IV Majoris Hebdomadæ inclusive:}

\lectioheading{Lectio brevis}{Is. 50, 6-7}

\lettrine[lines=2]{F}{ACIEM} meam non avérti ab increpántibus, et conspuéntibus in me. Dóminus Deus auxiliátor meus, et ídeo non sum confúsus. Tu autem.

\rubric{A Dominica Resurrectionis usque ad Vigiliam Ascensionis inclusive:}

\lectioheading{Lectio brevis}{Col. 3, 1-2}

\lettrine[lines=2]{S}{I} consurrexístis cum Christo, quæ sursum sunt quærite; ubi Christus est in déxtera Dei sedens: quæ sursum sunt sápite, non quæ super terram. Tu autem.

\rubric{In Officio autem cujuslibet Festi vel Octavæ, atque sanctæ Mariæ in Sabbato; itemque in Dominica infra Octavam Ascensionis, atque in Feria VI post Octavam Ascensionis, si de eis fiat Officium, et in Vigilia Pentecostes:}

\rubric{Pro Lectione brevi sumitur Capitulum, quod ad Nonam ponitur ipsius Officii currentis, et in fine dicitur: Tu autem, Dómine, ut supra.}

\rubric{Post Lectionem brevem dicitur:}

\V\ Adjutórium nostrum in nómine Dómini. \R\ Qui fecit cælum et terram.

\V\ Benedícite. \R\ Deus.

\benedlabel{Bened.} Dóminus nos benedícat, et ab omni malo deféndat, et ad vitam perdúcat ætérnam. Et fidélium ánimæ per misericórdiam Dei requiéscant in pace. \R\ Amen.

\rubric{Deinde dicitur tantum Pater noster secreto, nisi sequatur alia Hora.}

\hourheading{Ad Tertiam}

Pater noster et Ave.

\V\ Deus, in adjutórium.

\hymnusheading{Hymnus}

\lettrine[lines=2]{N}{UNC}, Sancte, nobis, Spíritus,\\
Unum Patri cum Fílio,\\
Dignáre promptus íngeri\\
Nostro refúsus péctori.

Os, lingua, mens, sensus, vigor,\\
Confessiónem pérsonent,\\
Flamméscat igne cáritas,\\
Accéndat ardor próximos.

Præsta, Pater piíssime,\\
Patríque compar Unice,\\
Cum Spíritu Paráclito\\
Regnans per omne sæculum. Amen.

\rubric{¶ In Festo Pentecostes et per Octavam, loco præcedentis, dicitur Hymnus} Veni, Creátor Spíritus, \rubric{ut in Proprio de Tempore.}

\rubric{Expleto Hymno, convéniens dicitur Antiphona usque ad Asteriscum *, et tres Psalmi, prouti Officium occurrens requirit; atque Psalmis absolutis, integra repetitur Antiphona.}

\variantheading{Tempore Quadragesimæ}

\rubric{In Officio dominicali, Antiphona et Capitulum ut in Proprio de Tempore.}

\rubric{In Officio autem feriali, dicitur quotidie sequens}

\antlabel{Ant.} Advenérunt nobis * dies pæniténtiæ, ad rediménda peccáta, ad salvándas ánimas.

\capitulumlabel{Capitulum}{Joel 2, 12-13}

\lettrine[lines=2]{C}{ONVERTIMINI} ad me in toto corde vestro, in jejúnio et fletu, et planctu. Et scíndite corda vestra, et non vestiménta vestra, ait Dóminus omnípotens.

\rubric{In Officio tamen dominicali et feriali dicitur sequens}

\R\ br. Ipse liberávit me * De láqueo venántium. Ipse. \V\ Et a verbo áspero. De. Glória Patri. Ipse.

\V\ Scápulis suis obumbrábit tibi. \R\ Et sub pennis ejus sperábis.

\variantheading{Tempore Passionis}

\rubric{In Officio dominicali, Antiphona et Capitulum ut in Proprio de Tempore.}

\rubric{In Officio feriali: In omnibus Feriis infra hebdomadam Passionis dicitur sequens}

\antlabel{Ant.} Judicásti, Dómine, * causam ánimæ meæ, defénsor vitæ meæ, Dómine, Deus meus.

\rubric{In Feria vero II, III et IV Majoris Hebdomadæ Antiphona secunda ex Laudibus Feriæ currentis ut in Proprio de Tempore.}

\rubric{In ipsis autem Feriis, infra utramque hebdomadam, dicitur quotidie sequens}

\capitulumlabel{Capitulum}{Jer. 17, 13}

\lettrine[lines=2]{D}{OMINE}, omnes qui te derelínquunt, confundéntur: recedéntes a te, in terra scribéntur: quóniam dereliquérunt venam aquárum vivéntium Dóminum.

\rubric{In Dominicis vero et in ómnibus Fériis, supérius indicátis, dicitur sequens}

\R\ br. Erue a frámea, * Deus, ánimam meam. Erue. \V\ Et de manu canis únicam meam. Deus. Erue. \V\ De ore leónis líbera me, Dómine. \R\ Et a córnibus unicórnium humilitátem meam.

\variantheading{Tempore Paschali}

\rubric{In Officio dominicali, a Dominica in Albis usque ad Dominicam V post Pascha inclusive, Capitulum ut in Proprio de Tempore.}

\rubric{In Officio feriali, a Feria II post Dominicam in Albis usque ad Vigiliam Ascensionis inclusive dicitur quotidie sequens}

\capitulumlabel{Capitulum}{Rom. 6, 9-10}

\lettrine[lines=2]{C}{HRISTUS} resúrgens ex mórtuis jam non móritur, mors illi ultra non dominábitur. Quod enim mórtuus est peccáto, mórtuus est semel: quod autem vivit, vivit Deo.

\rubric{In Officio tamen dominicali et feriali dicitur sequens}

\R\ br. Surréxit Dóminus de sepúlcro, * Allelúja, allelúja. Surréxit. \V\ Qui pro nobis pepéndit in ligno. Allelúja, allelúja. Glória Patri. Surréxit.

\V\ Surréxit Dóminus vere, allelúja. \R\ Et appáruit Simóni, allelúja.

\rubric{Expleto Responsorio brevi, si Preces non fuerint recitandæ, statim dicitur Versus} Dóminus vobíscum \rubric{cum Oratione et reliquis omnibus usque ad finem Horæ, ut infra.}

\rubric{Quando vero ad Laudes recitatæ sunt Preces, post Responsorium breve dicuntur sequentes flexis genibus:}

KYRIE, eléison. Christe, eléison. Kýrie, eléison.

Pater noster \rubric{secreto usque ad} \V\ Et ne nos indúcas in tentatiónem. \R\ Sed líbera nos a malo. \V\ Dómine, Deus virtútum, convérte nos. \R\ Et osténde fáciem tuam, et salvi érimus. \V\ Exsúrge, Christe, ádjuva nos. \R\ Et líbera nos propter nomen tuum. \V\ Dómine, exáudi oratiónem meam. \R\ Et clamor meus ad te véniat.

\rubric{Absolutis Precibus, vel, si Preces locum non habeant, post Responsorium breve statim dicitur:}

\V\ Dóminus vobíscum. \R\ Et cum spíritu tuo.

Orémus.

\rubric{Et dicitur Oratio conveniens.}

\rubric{Post Orationem additur:}

\V\ Dóminus vobíscum. \R\ Et cum spíritu tuo. \V\ Benedicámus Dómino. \R\ Deo grátias. \V\ Fidélium ánimæ per misericórdiam Dei requiéscant in pace. \R\ Amen.

\rubric{Deinde dicitur tantum Pater noster secreto, nisi sequatur alia Hora.}

\hourheading{Ad Sextam}

Pater noster et Ave.

\V\ Deus, in adjutórium.

\hymnusheading{Hymnus}

\lettrine[lines=2]{R}{ECTOR} potens, verax Deus,\\
Qui témperas rerum vices,\\
Splendóre mane illúminas,\\
Et ígnibus merídiem:

Exstíngue flammas lítium,\\
Aufer calórem nóxium,\\
Confer salútem córporum,\\
Verámque pacem córdium.

Præsta, Pater piíssime,\\
Patríque compar Unice,\\
Cum Spíritu Paráclito\\
Regnans per omne sæculum. Amen.

\rubric{Expleto Hymno, convéniens dicitur Antiphona usque ad Asteriscum *, et tres Psalmi, prouti Officium occurrens requirit; atque, Psalmis absolutis, integra repetitur Antiphona.}

\variantheading{Tempore Quadragesimæ}

\rubric{In Officio dominicali, Antiphona et Capitulum ut in Proprio de Tempore.}

\rubric{In Officio autem feriali, dicitur quotidie sequens}

\antlabel{Ant.} Commendémus nosmetípsos * in multa patiéntia, in jejúniis multis, per arma justítiæ.

\capitulumlabel{Capitulum}{Is. 55, 7}

\lettrine[lines=2]{D}{ERELINQUAT} ímpius viam suam, et vir iníquus cogitatiónes suas, et revertátur ad Dóminum, et miserébitur ejus, et ad Deum nostrum, quóniam multus est ad ignoscéndum.

\rubric{In Officio tamen dominicali et feriali dicitur sequens}

\R\ br. Scápulis suis * Obumbrábit tibi. Scápulis. \V\ Et sub pennis ejus sperábis. Obumbrábit tibi. Glória Patri. Scápulis.

\V\ Scuto circúmdabit te véritas ejus. \R\ Non timébis a timóre noctúrno.

\variantheading{Tempore Passionis}

\rubric{In Officio dominicali, Antiphona et Capitulum ut in Proprio de Tempore.}

\rubric{In Officio feriali: In omnibus Feriis infra hebdomadam Passionis dicitur sequens}

\antlabel{Ant.} Pópule meus, * quid feci tibi, aut quid moléstus fui? Respónde mihi.

\rubric{In Feria II, III et IV Majoris Hebdomadæ Antiphona tertia ex Laudibus Feriæ currentis ut in Proprio de Tempore.}

\rubric{In ipsis autem Feriis, infra utramque hebdomadam, dicitur quotidie sequens}

\capitulumlabel{Capitulum}{Jer. 17, 18}

\lettrine[lines=2]{C}{ONFUNDANTUR} qui me persequúntur, et non confúndar ego: páveant illi, et non páveam ego: induc super eos diem afflictiónis, et dúplici contritióne cóntere eos, Dómine, Deus noster.

\rubric{In Dominicis vero et in ómnibus Fériis, supérius indicátis, dicitur sequens}

\R\ br. De ore leónis * Líbera me, Dómine. De. \V\ Et a córnibus unicórnium humilitátem meam. Líbera. De. \V\ Ne perdas cum ímpiis, Deus, ánimam meam. \R\ Et cum viris sánguinum vitam meam.

\variantheading{Tempore Paschali}

\rubric{In Officio dominicali, a Dominica in Albis usque ad Dominicam V post Pascha inclusive, Capitulum ut in Proprio de Tempore.}

\rubric{In Officio feriali: a Feria II post Dominicam in Albis usque ad Vigiliam Ascensionis inclusive dicitur quotidie sequens}

\capitulumlabel{Capit.}{1 Cor. 15, 20-22}

\lettrine[lines=2]{C}{HRISTUS} resurréxit a mórtuis primítiæ dormiéntium: quóniam quidem per hóminem mors, et per hóminem resurréctio mortuórum. Et sicut in Adam omnes moriúntur, ita et in Christo omnes vivificabúntur.

\rubric{In Officio tamen dominicali et feriali dicitur sequens}

\R\ br. Surréxit Dóminus vere, * Allelúja, allelúja. Surréxit. \V\ Et appáruit Simóni. Allelúja, allelúja. Glória Patri. Surréxit.

\V\ Gavísi sunt discípuli, allelúja. \R\ Viso Dómino, allelúja.

\rubric{Expleto Responsorio brevi, si Preces non fuerint recitandæ, statim dicitur Versus} Dóminus vobíscum \rubric{cum Oratione et reliquis omnibus usque ad finem Horæ, ut infra.}

\rubric{Quando vero ad Laudes recitatæ sunt Preces, post Responsorium breve dicuntur sequentes flexis genibus:}

KYRIE, eléison. Christe, eléison. Kýrie, eléison.

Pater noster \rubric{secreto usque ad} \V\ Et ne nos indúcas in tentatiónem. \R\ Sed líbera nos a malo. \V\ Dómine, Deus virtútum, convérte nos. \R\ Et osténde fáciem tuam, et salvi érimus. \V\ Exsúrge, Christe, ádjuva nos. \R\ Et líbera nos propter nomen tuum. \V\ Dómine, exáudi oratiónem meam. \R\ Et clamor meus ad te véniat.

\rubric{Absolutis Precibus, vel, si Preces locum non habeant, post Responsorium breve statim dicitur:}

\V\ Dóminus vobíscum. \R\ Et cum spíritu tuo.

Orémus.

\rubric{Et dicitur Oratio conveniens.}

\rubric{Post Orationem additur:}

\V\ Dóminus vobíscum. \R\ Et cum spíritu tuo. \V\ Benedicámus Dómino. \R\ Deo grátias. \V\ Fidélium ánimæ per misericórdiam Dei requiéscant in pace. \R\ Amen.

\rubric{Deinde dicitur tantum Pater noster secreto, nisi sequatur alia Hora.}

\hourheading{Ad Nonam}

Pater noster et Ave.

\V\ Deus, in adjutórium.

\hymnusheading{Hymnus}

\lettrine[lines=2]{R}{ERUM}, Deus, tenax vigor,\\
Immótus in te pérmanens,\\
Lucis diúrnæ témpora\\
Successíbus detérminans:

Largíre lumen véspere,\\
Quo vita nusquam décidat,\\
Sed præmium mortis sacræ\\
Perénnis instet glória.

Præsta, Pater piíssime,\\
Patríque compar Unice,\\
Cum Spíritu Paráclito\\
Regnans per omne sæculum. Amen.

\rubric{Expleto Hymno, convéniens dicitur Antiphona usque ad Asteriscum *, et tres Psalmi, prouti Officium occurrens requirit; atque, Psalmis absolutis, integra repetitur Antiphona.}

\variantheading{Tempore Quadragesimæ}

\rubric{In Officio dominicali, Antiphona et Capitulum ut in Proprio de Tempore.}

\rubric{In Officio feriali: dicitur quotidie sequens}

\antlabel{Ant.} Per arma justítiæ * virtútis Dei commendémus nosmetípsos in multa patiéntia.

\capitulumlabel{Capitulum}{Is. 58, 7}

\lettrine[lines=2]{F}{RANGE} esuriénti panem tuum, et egénos vagósque induc in domum tuam: cum víderis nudum, óperi eum, et carnem tuam ne despéxeris.

\rubric{In Officio tamen dominicali et feriali dicitur sequens}

\R\ br. Scuto circúmdabit te * Véritas ejus. Scuto. \V\ Non timébis a timóre noctúrno. Véritas ejus. Glória Patri. Scuto.

\V\ Angelis suis Deus mandávit de te. \R\ Ut custódiant te in ómnibus viis tuis.

\variantheading{Tempore Passionis}

\rubric{In Officio dominicali, Antiphona et Capitulum ut in Proprio de Tempore.}

\rubric{In Officio feriali: In omnibus Feriis infra hebdomadam Passionis dicitur sequens}

\antlabel{Ant.} Numquid rédditur * pro bono malum, quia fodérunt fóveam ánimæ meæ?

\rubric{In Feria vero II, III et IV Majoris Hebdomadæ Antiphona quinta ex Laudibus Feriæ currentis ut in Proprio de Tempore.}

\rubric{In ipsis autem Feriis, infra utramque hebdomadam, dicitur quotidie sequens}

\capitulumlabel{Capitulum}{Jer. 18, 20}

\lettrine[lines=2]{R}{ECORDARE} quod stéterim in conspéctu tuo, ut lóquerer pro eis bonum, et avérterem indignatiónem tuam ab eis.

\rubric{In Dominicis vero et in ómnibus Fériis, supérius indicátis, dicitur sequens}

\R\ br. Ne perdas cum ímpiis, * Deus, ánimam meam. Ne. \V\ Et cum viris sánguinum vitam meam. Deus. Ne. \V\ Eripe me, Dómine, ab hómine malo. \R\ A viro iníquo éripe me.

\variantheading{Tempore Paschali}

\rubric{In Officio dominicali, a Dominica in Albis usque ad Dominicam V post Pascha inclusive, Capitulum ut in Proprio de Tempore.}

\rubric{In Officio feriali: a Feria II post Dominicam in Albis usque ad Vigiliam Ascensionis inclusive dicitur quotidie sequens}

\capitulumlabel{Capitulum}{1 Petri 3, 18}

\lettrine[lines=2]{C}{HRISTUS} semel pro peccátis nostris mórtuus est, justus pro injústis, ut nos offérret Deo, mortificátus quidem carne, vivificátus autem spíritu.

\rubric{In Officio tamen dominicali et feriali dicitur sequens}

\R\ br. Gavísi sunt discípuli, * Allelúja, allelúja. Gavísi. \V\ Viso Dómino. Allelúja, allelúja. Glória Patri. Gavísi.

\V\ Mane nobíscum, Dómine, allelúja. \R\ Quóniam advesperáscit, allelúja.

\rubric{Expleto Responsorio brevi, si Preces non fuerint recitandæ, statim dicitur Versus} Dóminus vobíscum \rubric{cum Oratione et reliquis omnibus usque ad finem Horæ, ut infra.}

\rubric{Quando vero ad Laudes recitatæ sunt Preces post Responsorium breve dicuntur sequentes flexis genibus:}

KYRIE, eléison. Christe, eléison. Kýrie, eléison.

Pater noster \rubric{secreto usque ad} \V\ Et ne nos indúcas in tentatiónem. \R\ Sed líbera nos a malo. \V\ Dómine, Deus virtútum, convérte nos. \R\ Et osténde fáciem tuam, et salvi érimus. \V\ Exsúrge, Christe, ádjuva nos. \R\ Et líbera nos propter nomen tuum. \V\ Dómine, exáudi oratiónem meam. \R\ Et clamor meus ad te véniat.

\rubric{Absolutis Precibus, vel, si Preces locum non habeant, post Responsorium breve statim dicitur:}

\V\ Dóminus vobíscum. \R\ Et cum spíritu tuo.

Orémus.

\rubric{Et dicitur Oratio conveniens.}

\rubric{Post Orationem additur:}

\V\ Dóminus vobíscum. \R\ Et cum spíritu tuo. \V\ Benedicámus Dómino. \R\ Deo grátias. \V\ Fidélium ánimæ per misericórdiam Dei requiéscant in pace. \R\ Amen.

\rubric{Deinde dicitur tantum Pater noster secreto, nisi sequatur alia Hora.}

\hourheading{Ad Vesperas}

Pater noster et Ave.

\V\ Deus, in adjutórium.

\rubric{Deinde, sub congruentibus Antiphonis, dicuntur quinque Psalmi, prouti Officium occurrens requirit.}

\rubric{Repetita vero post ultimum Psalmum Antiphona, dicuntur Capitulum, Hymnus et Versus, prouti Officium occurrens exigit.}

\variantheading{Tempore Quadragesimæ}

\rubric{In Officio dominicali, Capitulum ut in Proprio de Tempore.}

\rubric{In Officio feriali dicitur quotidie sequens}

\capitulumlabel{Capitulum}{Joel 2, 17}

\lettrine[lines=2]{I}{NTER} vestíbulum et altáre plorábunt sacerdótes, ministri Dómini, et dicent: Parce, Dómine, parce pópulo tuo: et ne des hereditátem tuam in oppróbrium, ut dominéntur eis natiónes.

\rubric{In Officio tamen dominicali et feriali dicitur sequens Hymnus et Versus:}

\hymnusheading{Hymnus}

\lettrine[lines=2]{A}{UDI}, benígne Cónditor,\\
Nostras preces cum flétibus,\\
In hoc sacro jejúnio\\
Fusas quadragenário.

Scrutátor alme córdium,\\
Infírma tu scis vírium:\\
Ad te reversis éxhibe\\
Remissiónis grátiam.

Multum quidem peccávimus,\\
Sed parce confiténtibus:\\
Ad nóminis laudem tui\\
Confer medélam lánguidis.

Concéde nostrum cónteri\\
Corpus per abstinéntiam;\\
Culpæ ut relínquant pábulum\\
Jejúna corda críminum.

\rubric{¶ Sequens Conclusio numquam mutatur:}

Præsta, beáta Trínitas,\\
Concéde, símplex Unitas;\\
Ut fructuósa sint tuis\\
Jejuniórum múnera. Amen.

\V\ Angelis suis Deus mandávit de te. \R\ Ut custódiant te in ómnibus viis tuis.

\variantheading{Tempore Passionis}

\rubric{In Officio dominicali, Capitulum ut in Proprio de Tempore.}

\rubric{In Officio feriali usque ad Feriam IV Majoris Hebdomadæ inclusive dicitur quotidie sequens}

\capitulumlabel{Capitulum}{Jer. 11, 20}

\lettrine[lines=2]{T}{U} autem, Dómine Sábaoth, qui júdicas juste, et probas renes et corda, vídeam ultiónem tuam ex eis: tibi enim revelávi causam meam, Dómine, Deus meus.

\rubric{In Officio tamen dominicali et feriali dicitur sequens Hymnus et Versus:}

\hymnusheading{Hymnus}

\lettrine[lines=2]{V}{EXILLA} Regis pródeunt:\\
Fulget Crucis mystérium,\\
Qua vita mortem pértulit,\\
Et morte vitam prótulit.

Quæ, vulneráta lánceæ\\
Mucróne diro, críminum\\
Ut nos laváret sórdibus,\\
Manávit unda et sánguine.

Impléta sunt quæ cóncinit\\
David fidéli cármine,\\
Dicéndo natiónibus:\\
Regnávit a ligno Deus.

Arbor decóra et fúlgida,\\
Ornáta Regis púrpura,\\
Elécta digno stípite\\
Tam sancta membra tángere.

Beáta, cujus bráchiis\\
Prétium pepéndit sæculi,\\
Statéra facta córporis,\\
Tulítque prædam tártari.

\rubric{¶ Sequens stropha dicitur flexis genibus; et última numquam mutátur.}

O Crux, ave, spes única,\\
Hoc Passiónis témpore\\
Piis adáuge grátiam,\\
Reísque dele crímina.

Te, fons salútis, Trínitas,\\
Collaúdet omnis spíritus:\\
Quibus Crucis victóriam\\
Largíris, adde præmium. Amen.

\V\ Eripe me, Dómine, ab hómine malo. \R\ A viro iníquo éripe me.

\variantheading{Tempore Paschali}

\rubric{In Officio dominicali a Dominica in Albis usque ad Dominicam V post Pascha inclusive, Capitulum ut in Proprio de Tempore.}

\rubric{In Officio feriali a Feria II post Dominicam in Albis usque ad Feriam III Rogationum inclusive dicitur quotidie sequens}

\capitulumlabel{Capitulum}{Rom. 6, 9-10}

\lettrine[lines=2]{C}{HRISTUS} resúrgens ex mórtuis jam non móritur, mors illi ultra non dominábitur. Quod enim mórtuus est peccáto, mórtuus est semel: quod autem vivit, vivit Deo.

\rubric{In Officio tamen dominicali et feriali dicitur sequens Hymnus et Versus:}

\hymnusheading{Hymnus}

\lettrine[lines=2]{A}{D} régias Agni dapes,\\
Stolis amícti cándidis,\\
Post tránsitum Maris Rubri,\\
Christo canámus Príncipi:

Divína cujus cáritas\\
Sacrum propínat sánguinem,\\
Almíque membra córporis\\
Amor sacérdos immolat.

Sparsum cruórem póstibus\\
Vastátor horret Angelus:\\
Fugítque divísum mare;\\
Mergúntur hostes flúctibus.

Jam Pascha nostrum Christus est,\\
Paschális idem víctima,\\
Et pura puris méntibus\\
Sinceritátis ázyma.

O vera cæli víctima,\\
Subjécta cui sunt tártara,\\
Solúta mortis víncula,\\
Recépta vitæ præmia.

Victor, subáctis ínferis,\\
Trophæa Christus éxplicat;\\
Cælóque apérto, súbditum\\
Regem tenebrárum trahit.

Ut sis perénne méntibus\\
Paschále, Jesu, gáudium,\\
A morte dira críminum,\\
Vitæ renátos líbera.

Deo Patri sit glória,\\
Et Fílio, qui a mórtuis\\
Surréxit, ac Paráclito,\\
In sempitérna sæcula. Amen.

\V\ Mane nobíscum, Dómine, allelúja. \R\ Quóniam advesperáscit, allelúja.

\rubric{Recitato Versu, dicitur cum Antiphona convenienti sequens}

\canticumheading{Canticum beatæ Mariæ Virginis}{Luc. 1, 46-55}

\lettrine[lines=2]{M}{AGNIFICAT} \textdagger\ ánima mea Dóminum:

Et exsultávit spíritus meus * in Deo, salutári meo.

Quia respéxit humilitátem ancíllæ suæ: * ecce enim, ex hoc beátam me dicent omnes generatiónes.

Quia fecit mihi magna qui potens est: * et sanctum nomen ejus.

Et misericórdia ejus a progénie in progénies * timéntibus eum.

Fecit poténtiam in bráchio suo: * dispérsit supérbos mente cordis sui.

Depósuit poténtes de sede, * et exaltávit húmiles.

Esuriéntes implévit bonis: * et dívites dimísit inánes.

Suscépit Israël, púerum suum, * recordátus misericórdiæ suæ.

Sicut locútus est ad patres nostros, * Abraham, et sémini ejus in sæcula.

Glória Patri.

\rubric{Expleto Cantico ac repetita Antiphona, si Preces non fuerint recitandæ, statim dicitur Versus} Dóminus vobíscum \rubric{cum Oratione et reliquis omnibus usque ad finem Horæ, ut infra.}

\rubric{In Feriis vero Quadragesimæ, et Passionis usque ad Feriam IV Majoris Hebdomadæ inclusive, si fiat Officium de Feria currenti, et quamvis in eo peragenda sit Commemoratio Duplicis aut Semiduplicis ad instar Simplicis redacti, post repetitam ad Magnificat Antiphonam dicuntur flexis genibus sequentes Preces, quæ alias omittuntur:}

KYRIE, eléison. Christe, eléison. Kýrie, eléison.

\rubric{Pater noster, quod dicitur a solo Hebdomadario totum clara voce, usque ad} \V\ Et ne nos indúcas in tentatiónem. \R\ Sed líbera nos a malo. \V\ Ego dixi: Dómine, miserére mei. \R\ Sana ánimam meam, quia peccávi tibi. \V\ Convértere, Dómine, úsquequo? \R\ Et deprecábilis esto super servos tuos. \V\ Fiat misericórdia tua, Dómine, super nos. \R\ Quemádmodum sperávimus in te. \V\ Sacerdótes tui induántur justítiam. \R\ Et sancti tui exsúltent.

\V\ Orémus pro beatíssimo Papa nostro N. \R\ Dóminus consérvet eum, et vivíficet eum, et beátum fáciat eum in terra, et non tradat eum in ánimam inimicórum ejus. \V\ Orémus et pro Antístite nostro N. \R\ Stet et pascat in fortitúdine tua, Dómine, in sublimitáte nóminis tui.

\rubric{¶ Romæ præcedens Versus cum suo Responsorio omittitur; alibi vero, ad litteram N., ab omnibus nomen diœcesani Episcopi exprimatur. Vacante Apostolica vel Episcopali Sede, alteruter vel uterque respectivus Versus cum suo Responsorio prætéritur.}

\V\ Dómine, salvum fac regem. \R\ Et exáudi nos in die, qua invocavérimus te. \V\ Salvum fac pópulum tuum, Dómine, et bénedic hereditáti tuæ. \R\ Et rege eos, et extólle illos usque in ætérnum. \V\ Meménto Congregatiónis tuæ. \R\ Quam possedísti ab inítio. \V\ Fiat pax in virtúte tua. \R\ Et abundántia in túrribus tuis. \V\ Orémus pro benefactóribus nostris. \R\ Retribúere dignáre, Dómine, ómnibus, nobis bona faciéntibus propter nomen tuum, vitam ætérnam. Amen. \V\ Orémus pro fidélibus defúnctis. \R\ Réquiem ætérnam dona eis, Dómine, et lux perpétua lúceat eis. \V\ Requiéscant in pace. \R\ Amen. \V\ Pro frátribus nostris abséntibus. \R\ Salvos fac servos tuos, Deus meus, sperántes in te. \V\ Pro afflíctis et captívis. \R\ Líbera eos, Deus Israël, ex ómnibus tribulatiónibus suis. \V\ Mitte eis, Dómine, auxílium de sancto. \R\ Et de Sion tuére eos. \V\ Dómine, Deus virtútum, convérte nos. \R\ Et osténde fáciem tuam, et salvi érimus. \V\ Exsúrge, Christe, ádjuva nos. \R\ Et líbera nos propter nomen tuum. \V\ Dómine, exáudi oratiónem meam. \R\ Et clamor meus ad te véniat.

\rubric{Absolutis Precibus, vel, si Preces locum non habeant, repetita Antiphona post Magnificat, statim dicitur:}

\V\ Dóminus vobíscum. \R\ Et cum spíritu tuo.

Orémus.

\rubric{Et dicitur Oratio conveniens.}

\rubric{Postea fiunt Commemorationes, si occurrant.}

\rubric{Deinde, extra Tempus Paschale, in Officio semiduplici aut simplici tam de Tempore quam de Sanctis, præter quam a Sabbato ante Dominicam I Passionis usque ad Feriam IV Majoris Hebdomadæ inclusive, ac præter omnes dies, in quibus habita sit Commemoratio Duplicis, etiam concurrentis, fit sequens}

\suffragiumheading{Suffragium de omnibus Sanctis}

\antlabel{Ant.} Beáta Dei Génitrix Virgo María, Sanctíque omnes intercédant pro nobis ad Dóminum.

\V\ Mirificávit Dóminus Sanctos suos. \R\ Et exaudívit eos clamántes ad se.

Orémus. \oratiolabel{Oratio}

\lettrine[lines=2]{A}{} \textsc{cunctis} nos, quæsumus, Dómine, mentis et córporis defénde perículis: et, intercedénte beáta et gloriósa semper Vírgine Dei Genitríce María, cum beáto Joseph, beátis Apóstolis tuis Petro et Paulo, atque beáto N. et ómnibus Sanctis, salútem nobis tríbue benígnus et pacem; ut, destrúctis adversitátibus et erróribus univérsis, Ecclésia tua secúra tibi sérviat libertáte. Per eúmdem Dóminum nostrum. \R\ Amen.

\rubric{¶ In hac Oratione, ad litteram N., exprimitur nomen Titularis propriæ Ecclesiæ, dummodo Titulus non sit Persona divina vel Mysterium Domini, aut de ipso ratione Festi factum non fuerit Officium vel Commemoratio, aut illius nomen non sit in eadem Suffragii Oratione expressum; ac nomina sanctorum Angelorum et sancti Joannis Baptistæ, si Titulares fuerint, præponuntur nomini sancti Joseph. In hisce autem omnibus casibus omittuntur verba atque beáto N.}

\rubric{Tempore autem Paschali, a Feria II post Dominicam in Albis usque ad Feriam III Rogationum inclusive, in Officio semiduplici aut simplici tam de Tempore quam de Sanctis, nisi habita sit Commemoratio Duplicis, etiam concurrentis, vel factum fuerit Officium aut Commemoratio cujusvis Octavæ, loco Suffragii, fit sequens}

\suffragiumheading{Commemoratio de Cruce}

\antlabel{Ant.} Crucifíxus surréxit a mórtuis, et redémit nos, allelúja, allelúja.

\V\ Dícite in natiónibus, allelúja. \R\ Quia Dóminus regnávit a ligno, allelúja.

Orémus. \oratiolabel{Oratio}

\lettrine[lines=2]{D}{EUS}, qui pro nobis Fílium tuum Crucis patíbulum subíre voluísti, ut inimíci a nobis expélleres potestátem: concéde nobis, fámulis tuis; ut resurrectiónis grátiam consequámur. Per eúmdem Dóminum. \R\ Amen.

\rubric{Post ultimam Orationem additur:}

\V\ Dóminus vobíscum. \R\ Et cum spíritu tuo. \V\ Benedicámus Dómino. \R\ Deo grátias. \V\ Fidélium ánimæ per misericórdiam Dei requiéscant in pace. \R\ Amen.

\rubric{Si post Vesperas immediate sequatur Completorium, dicto Versu} Fidélium ánimæ, \rubric{statim incipitur Versus} Jube, domne, benedícere, \rubric{ut infra ad Completorium; secus autem, si tunc terminetur Officium, dicitur tantum} Pater noster \rubric{secreto.}

\hourheading{Ad Completorium}

\rubric{Ante Completorium non dicitur} Pater noster et Ave María, \rubric{sicut in aliis Horis, sed absolute Lector incipit:}

\V\ Jube, domne, benedícere. \benedlabel{Bened.} Noctem quiétam, et finem perféctum concédat nobis Dóminus omnípotens. \R\ Amen.

\lectioheading{Lectio brevis}{1 Petr. 5, 8-9}

\lettrine[lines=2]{F}{RATRES}: Sóbrii estóte, et vigiláte: quia adversárius vester diábolus tamquam leo rúgiens círcuit, quærens quem dévoret: cui resístite fortes in fide. Tu autem, Dómine, miserére nobis. \R\ Deo grátias.

\V\ Adjutórium nostrum in nómine Dómini. \R\ Qui fecit cælum et terram.

Pater noster \rubric{quod dicitur totum secreto.}

\rubric{Deinde Hebdomadarius facit Confessionem:}

\lettrine[lines=2]{C}{ONFITEOR} Deo omnipoténti, beátæ Maríæ semper Vírgini, beáto Michaéli Archángelo, beáto Joánni Baptístæ, sanctis Apóstolis Petro et Paulo, ómnibus Sanctis, et vobis, fratres, quia peccávi nimis cogitatióne, verbo et ópere: mea culpa, mea culpa, mea máxima culpa. Ideo precor beátam Maríam semper Vírginem, beátum Michaélem Archángelum, beátum Joánnem Baptístam, sanctos Apóstolos Petrum et Paulum, omnes Sanctos, et vos, fratres, oráre pro me ad Dóminum, Deum nostrum.

\rubric{Chorus respondet:}

\lettrine[lines=2]{M}{ISEREATUR} tui omnípotens Deus, et, dimíssis peccátis tuis, perdúcat te ad vitam ætérnam. \R\ Amen.

\rubric{Deinde repetit Confessionem, et ubi dicitur:} vobis, fratres, et vos, fratres, \rubric{dicat:} tibi, pater, et te, pater.

\rubric{Facta Confessione a Choro Hebdomadarius dicit:}

\lettrine[lines=2]{M}{ISEREATUR} vestri omnípotens Deus, et, dimíssis peccátis vestris, perdúcat vos ad vitam ætérnam. \R\ Amen.

\lettrine[lines=2]{I}{NDULGENTIAM}, absolutiónem, et remissiónem peccatórum nostrórum tríbuat nobis omnípotens et miséricors Dóminus. \R\ Amen.

\rubric{Extra Chorum, si unus vel duo tantum recitent Officium, et in Choro Monialium, semel tantum ac simul ab omnibus ita fit Confessio:}

\lettrine[lines=2]{C}{ONFITEOR} Deo omnipoténti, beátæ Maríæ semper Vírgini, beáto Michaéli Archángelo, beáto Joánni Baptístæ, sanctis Apóstolis Petro et Paulo, et ómnibus Sanctis, quia peccávi nimis cogitatióne, verbo et ópere: mea culpa, mea culpa, mea máxima culpa. Ideo precor beátam Maríam semper Vírginem, beátum Michaélem Archángelum, beátum Joánnem Baptístam, sanctos Apóstolos Petrum et Paulum, et omnes Sanctos, oráre pro me ad Dóminum, Deum nostrum.

\rubric{Deinde dicitur:}

\lettrine[lines=2]{M}{ISEREATUR} nostri omnípotens Deus, et, dimíssis peccátis nostris, perdúcat nos ad vitam ætérnam. Amen.

\lettrine[lines=2]{I}{NDULGENTIAM}, absolutiónem, et remissiónem peccatórum nostrórum tríbuat nobis omnípotens et miséricors Dóminus. Amen.

\rubric{Facta Absolutione, dicitur:}

\V\ Convérte nos, Deus, salutáris noster. \R\ Et avérte iram tuam a nobis. \V\ Deus, in adjutórium.

\rubric{Deinde, inchoata convenienti Antiphona, dicuntur tres Psalmi, prouti Officium occurrens requirit.}

\rubric{Repetita integre post ultimum Psalmum Antiphona, dicitur sequens}

\hymnusheading{Hymnus}

\lettrine[lines=2]{T}{E} lucis ante términum,\\
Rerum Creátor, póscimus,\\
Ut pro tua cleméntia\\
Sis præsul et custódia.

Procul recédant sómnia,\\
Et nóctium phantásmata;\\
Hostémque nostrum cómprime,\\
Ne polluántur córpora.

Præsta, Pater piíssime,\\
Patríque compar Unice,\\
Cum Spíritu Paráclito\\
Regnans per omne sæculum. Amen.

\rubric{Absoluto Hymno, dicitur sequens Capitulum et Responsorium breve:}

\capitulumlabel{Capitulum}{Jer. 14, 9}

\lettrine[lines=2]{T}{U} autem in nobis es, Dómine, et nomen sanctum tuum invocátum est super nos: ne derelínquas nos, Dómine, Deus noster.

\R\ br. In manus tuas, Dómine, * Comméndo spíritum meum. In manus tuas, Dómine, comméndo spíritum meum. \V\ Redemísti nos, Dómine, Deus veritátis. Comméndo spíritum meum. Glória Patri, et Fílio, et Spirítui Sancto. In manus tuas, Dómine, comméndo spíritum meum.

\V\ Custódi nos, Dómine, ut pupíllam óculi. \R\ Sub umbra alárum tuárum prótege nos.

\rubric{A Sabbato ante Dominicam Passionis usque ad Feriam IV Majoris Hebdomadæ inclusive, si præcedentes Vesperæ fuerint de Dominica vel de Feria, in Responsorio brevi omittitur Glória Patri, et hujus loco repetitur Responsorium, ut sequitur:}

\R\ br. In manus tuas, Dómine, * Comméndo spíritum meum. In manus. \V\ Redemísti nos, Dómine, Deus veritátis. Comméndo. In manus.

\V\ Custódi nos, Dómine, ut pupíllam óculi. \R\ Sub umbra alárum tuárum prótege nos.

\rubric{A Sabbato ante Dominicam in Albis usque ad Feriam VI infra Octavam Pentecostes inclusive, dicitur sequens}

\R\ br. In manus tuas, Dómine, comméndo spíritum meum, * Allelúja, allelúja. In manus. \V\ Redemísti nos, Dómine, Deus veritátis. Allelúja, allelúja. Glória Patri. In manus.

\V\ Custódi nos, Dómine, ut pupíllam óculi, allelúja. \R\ Sub umbra alárum tuárum prótege nos, allelúja.

\rubric{Postea ad Canticum dicitur sequens}

\antlabel{Ant.} Salva nos.

\canticumheading{Canticum Simeonis}{Luc. 2, 29-32}

\lettrine[lines=2]{N}{UNC} dimíttis servum tuum, Dómine, * secúndum verbum tuum in pace:

Quia vidérunt óculi mei * salutáre tuum.

Quod parásti * ante fáciem ómnium populórum,

Lumen ad revelatiónem géntium, * et glóriam plebis tuæ Israël.

Glória Patri.

\antlabel{Ant.} Salva nos, Dómine, vigilántes, custódi nos dormiéntes; ut vigilémus cum Christo, et requiescámus in pace.

\rubric{Tempore Paschali in fine Antiphonæ additur} Allelúja.

\rubric{Deinde, si Preces non fuerint recitandæ, statim dicitur Versus} Dóminus vobíscum \rubric{cum subsequenti Oratione et reliquis omnibus usque ad finem Completorii, ut infra.}

\rubric{In quolibet Officio semiduplici, in Officio sanctæ Mariæ in Sabbato, in Festis simplicibus, atque in Feriis intra Hebdomadam communibus, post repetitam ad Nunc dimíttis Antiphonam dicuntur sequentes Preces; quæ tamen omittuntur in Feria VI post Octavam Ascensionis et omnibus insuper diebus, in quibus ad Vesperas habita sit Commemoratio Duplicis, vel factum fuerit Officium aut Commemoratio cujusvis Octavæ. Semper autem recitari debent, et flexis quidem genibus, in iis Feriis, in quibus ad Vesperas dictæ sint Preces feriales.}

KYRIE, eléison. Christe, eléison. Kýrie, eléison.

Pater noster \rubric{secreto usque ad} \V\ Et ne nos indúcas in tentatiónem. \R\ Sed líbera nos a malo.

Credo in Deum \rubric{secreto usque ad} \V\ Carnis resurrectiónem. \R\ Vitam ætérnam. Amen.

\V\ Benedíctus es, Dómine, Deus patrum nostrórum. \R\ Et laudábilis et gloriósus in sæcula. \V\ Benedicámus Patrem et Fílium cum Sancto Spíritu. \R\ Laudémus, et superexaltémus eum in sæcula. \V\ Benedíctus es, Dómine, in firmaménto cæli. \R\ Et laudábilis, et gloriósus, et superexaltátus in sæcula. \V\ Benedícat et custódiat nos omnípotens et miséricors Dóminus. \R\ Amen. \V\ Dignáre, Dómine, nocte ista. \R\ Sine peccáto nos custodíre. \V\ Miserére nostri, Dómine. \R\ Miserére nostri. \V\ Fiat misericórdia tua, Dómine, super nos. \R\ Quemádmodum sperávimus in te. \V\ Dómine, exáudi oratiónem meam. \R\ Et clamor meus ad te véniat.

\rubric{Absolutis Precibus, vel, si Preces locum non habeant, repetita Antiphona post Magnificat, statim dicitur:}

\V\ Dóminus vobíscum. \R\ Et cum spíritu tuo.

Orémus. \oratiolabel{Oratio}

\lettrine[lines=2]{V}{ISITA}, quæsumus, Dómine, habitatiónem istam, et omnes insídias inimíci ab ea longe repélle: Angeli tui sancti hábitent in ea, qui nos in pace custódiant; et benedíctio tua sit super nos semper. Per Dóminum nostrum. \R\ Amen.

\rubric{Post Orationem additur:}

\V\ Dóminus vobíscum. \R\ Et cum spíritu tuo. \V\ Benedicámus Dómino. \R\ Deo grátias.

\benedictiolabel{Benedictio} Benedícat et custódiat nos omnípotens et miséricors Dóminus, Pater, et Fílius, et Spíritus Sanctus. \R\ Amen.

\rubric{Et non dicitur Versus} Fidélium ánimæ, \rubric{sed immediate subjungitur una ex finalibus beátæ Mariæ Virginis Antiphonis, quæ infra pro diversitate Temporis assignantur.}

\rubric{Postea subjungitur:}

\V\ Divínum auxílium máneat semper nobíscum. \R\ Amen.

\rubric{Deinde dicitur secreto} Pater noster, Ave María et Credo.

\rubric{Si post Completorium immediate sequatur Matutinum, tunc bis dicitur} Pater noster, Ave María et Credo, \rubric{pro fine videlicet Officii diei præcedentis, et pro initio Officii diei sequentis.}

\sectionheading{Post singulas Horas}

\rubric{In fine Laudum, vel, si post Laudes immediate subsequatur alia Hora, in fine ultimæ Horæ, sicuti et semper in fine Completorii, quamvis immediate sequatur Matutinum diei sequentis, tam in publica quam in privata divini Officii recitatione; atque in publica insuper recitatione post alias quascumque Horas, si tunc terminetur Officium et discedatur a Choro, dicitur sequens Antiphona finalis beátæ Mariæ Virginis.}

\rubric{I Post Laudes ac ceteras Horas, præter Completorium, dicto Versu} Fidélium ánimæ \rubric{et Oratione Dominica secreto recitata, subjungitur:}

\V\ Dóminus det nobis suam pacem. \R\ Et vitam ætérnam. Amen.

\rubric{Atque statim dicitur Antiphona finalis, ut infra.}

\rubric{II Post Completorium vero, recitata Benedictione} Benedícat et custódiat nos, \rubric{immediate subjungitur Antiphona finalis, ut infra.}

\sectionheading{Antiphonæ finales Beatæ Mariæ Virginis}

\rubric{1 A Vesperis Sabbati ante Dominicam I Quadragesimæ usque ad Completorium Feriæ IV Majoris Hebdomadæ inclusive:}

\antiphonaheading{Antiphona}

\lettrine[lines=2]{A}{VE}, Regína cælórum,\\
Ave, Dómina Angelórum:\\
Salve, radix, salve, porta,\\
Ex qua mundo lux est orta:\\
Gaude, Virgo gloriósa,\\
Super omnes speciósa,\\
Vale, o valde decóra,\\
Et pro nobis Christum exóra.

\V\ Dignáre me laudáre te, Virgo sacráta. \R\ Da mihi virtútem contra hostes tuos.

Orémus. \oratiolabel{Oratio}

\lettrine[lines=2]{C}{ONCEDE}, miséricors Deus, fragilitáti nostræ præsídium: ut, qui sanctæ Dei Genitrícis memóriam ágimus; intercessiónis ejus auxílio, a nostris iniquitátibus resurgámus. Per eúmdem Christum Dóminum nostrum. \R\ Amen.

\rubric{2 A Completorio Sabbati sancti usque ad Nonam Sabbati infra Octavam Pentecostes inclusive:}

\antiphonaheading{Antiphona}

\lettrine[lines=2]{R}{EGINA} cæli, lætáre, allelúja;\\
Quia quem meruísti portáre, allelúja,\\
Resurréxit, sicut dixit, allelúja:\\
Ora pro nobis Deum, allelúja.

\V\ Gaude et lætáre, Virgo María, allelúja. \R\ Quia surréxit Dóminus vere, allelúja.

Orémus. \oratiolabel{Oratio}

\lettrine[lines=2]{D}{EUS}, qui per resurrectiónem Fílii tui, Dómini nostri Jesu Christi, mundum lætificáre dignátus es: præsta, quæsumus; ut, per ejus Genitrícem Vírginem Maríam, perpétuæ capiámus gáudia vitæ. Per eúmdem Christum, Dóminum nostrum. \R\ Amen.

\rubric{Absoluta vero Antiphona finali cum suis Versu et Oratione, concluditur:}

\V\ Divínum auxílium máneat semper nobíscum. \R\ Amen.

\sectionheading{Post Divinum Officium}

\rubric{Cum expletum fuerit Officium, laudabiliter dicitur sequens Oratio; pro qua Summus Pontifex Pius Papa XI sub die 1 Decembris 1933 indulgentiam trium annorum concessit. Præterea Clericis in sacris constitutis et Sacerdotibus eam devote post Officium recitantibus Summus Pontifex Leo X defectus et culpas in eo persolvendo ex humana fragilitate contractas, indulsit. Dicitur autem flexis semper genibus in privata etiam recitatione, præter quam ab iis, qui ob certam infirmitatem vel gravioris impedimenti causam nequeant genuflectere.}

\oratioheading{Oratio}

\lettrine[lines=2]{S}{ACROSANCTAE} et individuæ Trinitáti, crucifíxi Dómini nostri Jesu Christi humanitáti, beatíssimæ et gloriosíssimæ semperque Vírginis Maríæ fecúndæ integritáti, et ómnium Sanctórum universitáti sit sempitérna laus, honor, virtus et glória ab omni creatúra, nobísque remíssio ómnium peccatórum, per infiníta sæcula sæculórum. \R\ Amen.

\V\ Beáta víscera Maríæ Vírginis, quæ portavérunt ætérni Patris Fílium. \R\ Et beáta úbera, quæ lactavérunt Christum Dóminum.

\rubric{Et dicitur secreto} Pater noster et Ave María.

% [ornamental image ending Ordinarium Divini Officii]

% =====================================================================
% END OF BATCH 1 -- Ordinarium Divini Officii (verna.md lines 1-1745)
%
% Notes:
%  - verna.md's Ordinarium uses accented rubric text throughout, unlike
%    hiemalis.md's largely unaccented Ordinarium rubrics; every word in
%    this batch was transcribed fresh from verna.md, not copied from
%    hiemalis.tex, per the confirmed accentuation difference between
%    the two source files.
%  - Corrected one evident OCR artifact: verna.md line 7 has a Greek
%    small letter omicron (U+03BF) in "adservátο" where Latin "o" is
%    required ("adserváto"); silently normalized here.
%  - verna.md's Ordinarium contains additional Lenten/Passiontide/
%    Paschal-season material (Invitatoria, hymns, capitula, responsoria
%    brevia, the Commemoratio de Cruce) not present in hiemalis.md's
%    Ordinarium, since Pars Verna covers Lent through Pentecost. All
%    such material is new transcription, not adapted from hiemalis.tex.
%  - No [REF: p. N] tags or bare bracketed page numbers requiring
%    resolution appeared in this batch, except one bare "[3]" marker
%    after each seasonal Invitatorium/Ps. 94 Veníte line (verna.md
%    lines 89, 144, 198), which is not a cross-reference but a psalm-
%    tone/repetition-count marker from the source; preserved literally.
%  - verna.md's Ordinarium has no <!-- p. N --> page-marker comments or
%    printer's-signature marks (unlike hiemalis.md), so no \label{p:N}
%    or \sectionornament page-break markers were inserted in this
%    batch; the two [IMAGE] markers bracketing the Ordinarium (verna.md
%    lines 1 and 1744) are rendered as plain LaTeX comments.
%  - Lettrine splits follow the same first-letter-only convention
%    established in hiemalis.tex (e.g. V/ENITE, C/HRISTUS, C/ONFITEOR),
%    cross-checked against hiemalis.tex precedent where the same
%    incipit recurs; single-letter incipits ("A VINCULIS") use the
%    lettrine{X}{} + \textsc{} pattern per hiemalis.tex precedent, but
%    "O Sol salútis" and "O Crux, ave" were rendered with plain text
%    after the lettrine (no \textsc{}) because only the single-letter
%    word itself is capitalized in verna.md's source, not the whole
%    following word.
%  - Two commands not previously catalogued in the summary legend but
%    confirmed present in hiemalis.tex and reused here: \absolutioheading{}
%    for the standalone "Absolutio" heading (distinct from
%    \absolutiolabel{} used for abbreviated inline repeats),
%    \oratiolabel{} for the inline "Orémus. Oratio" pattern (distinct
%    from \oratioheading{} for a standalone "Oratio" heading line), and
%    \benedictiolabel{} for the full-word "Benedictio" label (distinct
%    from \benedlabel{} used for the abbreviated "Bened." label).
%
% NEXT BATCH (Batch 2) will begin at verna.md line 1746:
% "# Psalterium Breviarii Romani", "## Dominica".
% =====================================================================

% [ornamental image beginning next part: Psalterium Breviarii Romani]

% ===================================================================
% BATCH 2: Psalterium Breviarii Romani, Dominica -- Ad Matutinum
% (verna.md lines 1746-2236)
% ===================================================================

\partheading{Psalterium Breviarii Romani}
\partsubtitle{Per omnes ac singulos hebdomadæ dies dispositum}

\feriaheading{Dominica}

\hourheading{Ad Matutinum}

Pater, Ave et Credo.

\V\ Dómine, lábia mea.

\V\ Deus, in adjutórium.

\variantheading{Tempore Quadragesimæ}

\invitatlabel{Invit.} Non sit vobis vanum mane súrgere ante lucem: * Quia promísit Dóminus corónam vigilántibus. \psalmlabel{Ps. 94} Veníte. [3]

\hymnusheading{Hymnus}

\lettrine[lines=2]{E}{X} more docti mýstico\\
Servémus hoc jejúnium,\\
Deno diérum círculo\\
Ducto quater notíssimo.

Lex et prophétæ prímitus\\
Hoc prætulérunt, póstmodum\\
Christus sacrávit, ómnium\\
Rex atque factor témporum.

Utámur ergo párcius\\
Verbis, cibis et pótibus,\\
Somno, jocis, et árctius\\
Perstémus in custódia.

Vitémus autem nóxia,\\
Quæ súbruunt mentes vagas:\\
Nullúmque demus cállidi\\
Hostis locum tyránnidi.

Flectámus iram víndicem,\\
Plorémus ante Júdicem,\\
Clamémus ore súpplici,\\
Dicámus omnes cérnui:

Nostris malis offéndimus\\
Tuam, Deus, cleméntiam:\\
Effúnde nobis désuper,\\
Remíssor, indulgéntiam.

Meménto quod sumus tui,\\
Licet cadúci, plásmatis:\\
Ne des honórem nóminis\\
Tui, precámur, álteri.

Laxa malum, quod fécimus,\\
Auge bonum, quod póscimus:\\
Placére quo tandem tibi\\
Possímus hic, et pérpetim.

\rubric{¶ Sequens Conclusio numquam mutatur.}

Præsta, beáta Trínitas,\\
Concéde, símplex Unitas,\\
Ut fructuósa sint tuis\\
Jejuniórum múnera. Amen.

\variantheading{Tempore Passionis}

\invitatlabel{Invit.} Hódie, si vocem Dómini audiéritis, * Nolíte obduráre corda vestra. \psalmlabel{Ps. 94} Veníte. [3]

\rubric{In quarto Psalmi versu non dicitur: Hódie, si vocem ejus audiéritis, nolíte obduráre corda vestra; sed, post repetitum a Choro Invitatorium, statim subjúngitur: Sicut in exacerbatióne.}

\rubric{In fine autem hujus Psalmi non dicitur} Glória Patri, \rubric{quando fit Officium de Dominica; sed repetitur Invitatorium, ut alias, in fine.}

\hymnusheading{Hymnus}

\lettrine[lines=2]{P}{ANGE}, lingua, gloriósi\\
Láuream certáminis,\\
Et super Crucis trophǽo\\
Dic triúmphum nóbilem:\\
Quáliter Redémptor orbis\\
Immolátus vícerit.

De paréntis protoplásti\\
Fraude Factor condolens,\\
Quando pomi noxiális\\
In necem morsu ruit:\\
Ipse lignum tunc notávit,\\
Damna ligni ut sólveret.

Hoc opus nostræ salútis\\
Ordo depopóscerat;\\
Multifórmis proditóris\\
Ars ut artem fálleret,\\
Et medélam ferret inde,\\
Hostis unde læserat.

Quando venit ergo sacri\\
Plenitúdo témporis,\\
Missus est ab arce Patris\\
Natus, orbis Cónditor;\\
Atque ventre virgináli\\
Carne amíctus pródiit.

Vagit infans inter arcta\\
Cónditus præsépia:\\
Membra pannis involúta\\
Virgo Mater álligat:\\
Et Dei manus pedésque\\
Stricta cingit fáscia.

Sempitérna sit beátæ\\
Trinitáti glória,\\
Æqua Patri, Filióque;\\
Par decus Paráclito:\\
Uníus Triníque nomen\\
Laudet univérsitas. Amen.

\variantheading{Tempore Paschali}

\invitatlabel{Invit.} Surréxit Dóminus vere, * Allelúja. \psalmlabel{Ps. 94} Veníte. [3]

\hymnusheading{Hymnus}

\lettrine[lines=2]{R}{EX} sempitérne cælitum,\\
Rerum Creátor ómnium,\\
Æquális ante sæcula\\
Semper Parénti Fílius:

Nascénte qui mundo faber\\
Imáginem vultus tui\\
Tradens Adámo, nóbilem\\
Limo jugásti spíritum.

Cum livor et fraus dæmonis\\
Fœdásset humánum genus:\\
Tu, carne amíctus, pérditam\\
Formam refórmas ártifex.

Qui, natus olim e Vírgine,\\
Nunc e sepúlcro násceris,\\
Tecúmque nos a mórtuis\\
Jubes sepúltos súrgere.

Qui, pastor ætérnus, gregem\\
Aqua lavas baptísmatis:\\
Hæc est lavácrum méntium:\\
Hæc est sepúlcrum críminum.

Nobis diu qui débitæ\\
Redémptor affíxus cruci,\\
Nostræ dedísti pródigus\\
Prétium salútis sánguinem.

Ut sis perénne méntibus\\
Paschále, Jesu, gáudium,\\
A morte dira críminum\\
Vitæ renátos líbera.

Deo Patri sit glória,\\
Et Fílio, qui a mórtuis\\
Surréxit, ac Paráclito,\\
In sempitérna sæcula. Amen.

\nocturnoheading{In I Nocturno}

\rubric{Extra Tempus Paschale}

\antlabel{Ant.} Beátus vir. \textdagger

\rubric{¶ Et non repetitur in Psalmo; quod semper fit, quando Antiphona incipit a primis Psalmi vel Cantici verbis. Et Psalmus vel Canticum tunc incipitur ab eo verbo, ante quod vel inchoata vel integra desinit Antiphona; si tamen eadem sint verba, et continuetur Antiphona cum Psalmo vel Cantico, neque additur in fine ipsius Antiphonæ Allelúja.}

\variantheading{Tempore Paschali}

\antlabel{Ant.} Allelúja, * lapis revolútus est, allelúja: ab óstio monuménti, allelúja, allelúja.

\rubric{Sub una tantum Antiphona dicuntur tres Psalmi in unoquoque Nocturno, etiam in Feriis.}

\rubric{Antiphonæ autem cum Psalmis et Versibus, tam intra quam extra Tempus Paschale assignatæ, adhibentur semper, extra Octavam Ascensionis, quando fit Officium de Dominica seu minori seu majori.}

\rubric{¶ Asteriscus *, in Psalmorum versibus et in aliis Officii partibus insertus, denotat pausam cantus et recitationis in Choro; præter quam in Antiphonis, in quibus nempe significat tum præintonationem Antiphonæ, tum Antiphonam ipsam ante Psalmum seu Canticum eousque producendam esse ad omnes et singulas Horas in Officio ritus semiduplicis ac simplicis, atque ad Horas tantum minores in Officio ritus duplicis.}

\psalmheading{Psalmus 1}

\lettrine[lines=2]{B}{EATUS} vir, \textdagger\ qui non ábiit in consílio impiórum, et in via peccatórum non stetit, * et in cáthedra pestiléntiæ non sedit:

Sed in lege Dómini volúntas ejus, * et in lege ejus meditábitur die ac nocte.

Et erit tamquam lignum, quod plantátum est secus decúrsus aquárum, * quod fructum suum dabit in témpore suo:

Et fólium ejus non défluet: * et ómnia quæcúmque fáciet, prosperabúntur.

Non sic ímpii, non sic: * sed tamquam pulvis, quem prójicit ventus a fácie terræ.

Ideo non resúrgent ímpii in judício: * neque peccatóres in concílio justórum.

Quóniam novit Dóminus viam justórum: * et iter impiórum períbit.

Glória Patri, et Fílio, * et Spirítui Sancto.

Sicut erat in princípio, et nunc, et semper, * et in sæcula sæculórum. Amen.

\rubric{¶ Hic versus Glória Patri semper dicitur in fine omnium Psalmorum et Canticorum, nisi aliter notetur.}

\antlabel{Ant.} Beátus vir, qui in lege Dómini meditátur.

\antlabel{Ant.} Servíte Dómino.

\psalmheading{Psalmus 2}

\lettrine[lines=2]{Q}{UARE} fremuérunt gentes, * et pópuli meditáti sunt inánia?

Astitérunt reges terræ, et príncipes convenérunt in unum * advérsus Dóminum, et advérsus Christum ejus.

Dirumpámus víncula eórum: * et projiciámus a nobis jugum ipsórum.

Qui hábitat in cælis, irridébit eos: * et Dóminus subsannábit eos.

Tunc loquétur ad eos in ira sua, * et in furóre suo conturbábit eos.

Ego autem constitútus sum Rex ab eo super Sion montem sanctum ejus, * prædicans præcéptum ejus.

Dóminus dixit ad me: * Fílius meus es tu, ego hódie génui te.

Póstula a me, et dabo tibi gentes hereditátem tuam, * et possessiónem tuam términos terræ.

Reges eos in virga férrea, * et tamquam vas fíguli confrínges eos.

Et nunc, reges, intellígite: * erudímini, qui judicátis terram.

Servíte Dómino in timóre: * et exsultáte ei cum tremóre.

Apprehéndite disciplínam, nequándo irascátur Dóminus, * et pereátis de via justa.

Cum exárserit in brevi ira ejus, * beáti omnes qui confídunt in eo.

\antlabel{Ant.} Servíte Dómino in timóre: et exsultáte ei cum tremóre.

\antlabel{Ant.} Exsúrge.

\psalmheading{Psalmus 3}

\lettrine[lines=2]{D}{OMINE}, quid multiplicáti sunt qui tríbulant me? * multi insúrgunt advérsum me.

Multi dicunt ánimæ meæ: * Non est salus ipsi in Deo ejus.

Tu autem, Dómine, suscéptor meus es, * glória mea, et exáltans caput meum.

Voce mea ad Dóminum clamávi: * et exaudívit me de monte sancto suo.

Ego dormívi, et soporátus sum: * et exsurréxi, quia Dóminus suscépit me.

Non timébo míllia pópuli circumdántis me: * exsúrge, Dómine, salvum me fac, Deus meus.

Quóniam tu percussísti omnes adversántes mihi sine causa: * dentes peccatórum contrivísti.

Dómini est salus: * et super pópulum tuum benedíctio tua.

\rubric{Extra Tempus Paschale}

\antlabel{Ant.} Exsúrge, Dómine, salvum me fac, Deus meus.

\rubric{In Quadragesima}

\V\ Ipse liberávit me de láqueo venántium. \R\ Et a verbo áspero.

\rubric{Tempore Passionis}

\V\ Erue a frámea, Deus, ánimam meam. \R\ Et de manu canis únicam meam.

\rubric{Tempore Paschali}

\antlabel{Ant.} Allelúja, * lapis revolútus est, allelúja: ab óstio monuménti, allelúja, allelúja.

\V\ Surréxit Dóminus de sepúlcro, allelúja. \R\ Qui pro nobis pepéndit in ligno, allelúja.

\absolutiolabel{Absol.} Exáudi, Dómine.

\benedlabel{Bened.} 1 Benedictióne perpétua. 2 Unigénitus Dei Fílius. 3 Spíritus Sancti grátia.

\nocturnoheading{In II Nocturno}

\rubric{Extra Tempus Paschale}

\antlabel{Ant.} Quam admirábile.

\rubric{Tempore Paschali}

\antlabel{Ant.} Allelúja, * quem quæris, múlier? allelúja: vivéntem cum mórtuis, allelúja, allelúja.

\psalmheading{Psalmus 8}

\lettrine[lines=2]{D}{OMINE}, Dóminus noster, * quam admirábile est nomen tuum in univérsa terra!

Quóniam eleváta est magnificéntia tua, * super cælos.

Ex ore infántium et lacténtium perfecísti laudem propter inimícos tuos, * ut déstruas inimícum et ultórem.

Quóniam vidébo cælos tuos, ópera digitórum tuórum: * lunam et stellas, quæ tu fundásti.

Quid est homo, quod memor es ejus? * aut fílius hóminis, quóniam vísitas eum?

Minuísti eum paulo minus ab Angelis, glória et honóre coronásti eum: * et constituísti eum super ópera mánuum tuárum.

Omnia subjecísti sub pédibus ejus, * oves et boves univérsas: ínsuper et pécora campi.

Vólucres cæli, et pisces maris, * qui perámbulant sémitas maris.

Dómine, Dóminus noster, * quam admirábile est nomen tuum in univérsa terra!

\antlabel{Ant.} Quam admirábile est nomen tuum, Dómine, in univérsa terra!

\antlabel{Ant.} Sedísti super thronum.

\psalmheading{Psalmus 9, i}

\lettrine[lines=2]{C}{ONFITEBOR} tibi, Dómine, in toto corde meo: * narrábo ómnia mirabília tua.

Lætábor et exsultábo in te: * psallam nómini tuo, Altíssime,

In converténdo inimícum meum retrórsum: * infirmabúntur, et períbunt a fácie tua.

Quóniam fecísti judícium meum et causam meam: * sedísti super thronum, qui júdicas justítiam.

Increpásti gentes, et périit ímpius: * nomen eórum delésti in ætérnum, et in sæculum sæculi.

Inimíci defecérunt frámeæ in finem: * et civitátes eórum destruxísti.

Périit memória eórum cum sónitu: * et Dóminus in ætérnum pérmanet.

Parávit in judício thronum suum: * et ipse judicábit orbem terræ in æquitáte, judicábit pópulos in justítia.

Et factus est Dóminus refúgium páuperi: * adjútor in opportunitátibus, in tribulatióne.

Et sperent in te qui novérunt nomen tuum: * quóniam non dereliquísti quæréntes te, Dómine.

\antlabel{Ant.} Sedísti super thronum, qui júdicas justítiam.

\antlabel{Ant.} Exsúrge, Dómine.

\psalmheading{Psalmus 9, ii}

\lettrine[lines=2]{P}{SALLITE} Dómino, qui hábitat in Sion: * annuntiáte inter gentes stúdia ejus:

Quóniam requírens sánguinem eórum recordátus est: * non est oblítus clamórem páuperum.

Miserére mei, Dómine: * vide humilitátem meam de inimícis meis.

Qui exáltas me de portis mortis, * ut annúntiem omnes laudatiónes tuas in portis fíliæ Sion.

Exsultábo in salutári tuo: * infíxæ sunt gentes in intéritu, quem fecérunt.

In láqueo isto, quem abscondérunt, * comprehénsus est pes eórum.

Cognoscétur Dóminus judícia fáciens: * in opéribus mánuum suárum comprehénsus est peccátor.

Convertántur peccatóres in inférnum, * omnes gentes quæ obliviscúntur Deum.

Quóniam non in finem oblívio erit páuperis: * patiéntia páuperum non períbit in finem.

Exsúrge, Dómine, non confortétur homo: * judicéntur gentes in conspéctu tuo.

Constítue, Dómine, legislatórem super eos: * ut sciant gentes quóniam hómines sunt.

\rubric{Extra Tempus Paschale}

\antlabel{Ant.} Exsúrge, Dómine, non prævaleat homo.

\rubric{In Quadragesima}

\V\ Scápulis suis obumbrábit tibi. \R\ Et sub pennis ejus sperábis.

\rubric{Tempore Passionis}

\V\ De ore leónis líbera me, Dómine. \R\ Et a córnibus unicórnium humilitátem meam.

\rubric{Tempore Paschali}

\antlabel{Ant.} Allelúja, quem quæris, múlier? allelúja: vivéntem cum mórtuis, allelúja, allelúja.

\V\ Surréxit Dóminus vere, allelúja. \R\ Et appáruit Simóni, allelúja.

\absolutiolabel{Absol.} Ipsíus píetas et.

\benedlabel{Bened.} 4 Deus Pater omnípotens. 5 Christus perpétuæ. 6 Ignem sui amóris.

\nocturnoheading{In III Nocturno}

\rubric{Extra Tempus Paschale}

\antlabel{Ant.} Ut quid, Dómine. \textdagger

\rubric{Tempore Paschali}

\antlabel{Ant.} Allelúja, * noli flere, María, allelúja: resurréxit Dóminus, allelúja, allelúja.

\psalmheading{Psalmus 9, iii}

\lettrine[lines=2]{U}{T} quid, Dómine, \textdagger\ recessísti longe, * déspicis in opportunitátibus, in tribulatióne?

Dum supérbit ímpius, incénditur pauper: * comprehendúntur in consíliis quibus cógitant.

Quóniam laudátur peccátor in desidériis ánimæ suæ: * et iníquus benedícitur.

Exacerbávit Dóminum peccátor, * secúndum multitúdinem iræ suæ non quæret.

Non est Deus in conspéctu ejus: * inquinátæ sunt viæ illíus in omni témpore.

Auferúntur judícia tua a fácie ejus: * ómnium inimicórum suórum dominábitur.

Dixit enim in corde suo: * Non movébor a generatióne in generatiónem sine malo.

Cujus maledictióne os plenum est, et amaritúdine, et dolo: * sub lingua ejus labor et dolor.

Sedet in insídiis cum divítibus in occúltis, * ut interfíciat innocéntem.

Oculi ejus in páuperem respíciunt: * insidiátur in abscóndito, quasi leo in spelúnca sua.

Insidiátur ut rápiat páuperem: * rápere páuperem, dum áttrahit eum.

In láqueo suo humiliábit eum: * inclinábit se, et cadet, cum dominátus fúerit páuperum.

Dixit enim in corde suo: * Oblítus est Deus, avértit fáciem suam ne vídeat in finem.

\antlabel{Ant.} Ut quid, Dómine, recessísti longe?

\antlabel{Ant.} Exsúrge. \textdagger

\psalmheading{Psalmus 9, iv}

\lettrine[lines=2]{E}{XSURGE}, \textdagger\ Dómine Deus, exaltétur manus tua: * ne obliviscáris páuperum.

Propter quid irritávit ímpius Deum? * dixit enim in corde suo: Non requíret.

Vides quóniam tu labórem et dolórem consíderas: * ut tradas eos in manus tuas.

Tibi derelíctus est pauper: * órphano tu eris adjútor.

Cóntere bráchium peccatóris et malígni: * quærétur peccátum illíus, et non inveniétur.

Dóminus regnábit in ætérnum, et in sæculum sæculi: * períbitis, gentes, de terra illíus.

Desidérium páuperum exaudívit Dóminus: * præparatiónem cordis eórum audívit auris tua.

Judicáre pupíllo et húmili, * ut non appónat ultra magnificáre se homo super terram.

\antlabel{Ant.} Exsúrge, Dómine Deus, exaltétur manus tua.

\antlabel{Ant.} Justus Dóminus.

\psalmheading{Psalmus 10}

\lettrine[lines=2]{I}{N} Dómino confído: quómodo dícitis ánimæ meæ: * Tránsmigra in montem sicut passer?

Quóniam ecce peccatóres intendérunt arcum, paravérunt sagíttas suas in pháretra, * ut sagíttent in obscúro rectos corde.

Quóniam quæ perfecísti, destruxérunt: * justus autem quid fecit?

Dóminus in templo sancto suo, * Dóminus in cælo sedes ejus:

Oculi ejus in páuperem respíciunt: * pálpebræ ejus intérrogant fílios hóminum.

Dóminus intérrogat justum et ímpium: * qui autem díligit iniquitátem, odit ánimam suam.

Pluet super peccatóres láqueos: * ignis, et sulphur, et spíritus procellárum pars cálicis eórum.

Quóniam justus Dóminus, et justítias diléxit: * æquitátem vidit vultus ejus.

\rubric{Extra Tempus Paschale}

\antlabel{Ant.} Justus Dóminus, et justítiam diléxit.

\rubric{In Quadragesima}

\V\ Scuto circúmdabit te véritas ejus. \R\ Non timébis a timóre noctúrno.

\rubric{Tempore Passionis}

\V\ Ne perdas cum ímpiis, Deus, ánimam meam. \R\ Et cum viris sánguinum vitam meam.

\rubric{Tempore Paschali}

\antlabel{Ant.} Allelúja, noli flere, María, allelúja: resurréxit Dóminus, allelúja, allelúja.

\V\ Gavísi sunt discípuli, allelúja. \R\ Viso Dómino, allelúja.

\absolutiolabel{Absolutio} A vínculis.

\benedlabel{Bened.} 7 Evangélica léctio. 8 Divínum auxílium. 9 Ad societátem cívium.

% ===================================================================
% BATCH 3: Psalterium, Dominica -- Ad Laudes I and II
% (verna.md lines 2237-2602)
% ===================================================================

\hourheading{Ad Laudes I}

\textit{In Dominicis Temporis Paschalis}

(Pater noster \rubric{et} Ave).

\V\ Deus, in adjutórium.

\antlabel{Ant.} Allelúja, * allelúja, allelúja.

\rubric{Sub quæ sola Antiphona Allelúja dicuntur omnes Psalmi cum Cantico usque ad Capitulum, etiam in Feriis.}

\rubric{¶ Antiphona cum Psalmis assignata adhibetur semper extra Octavam Ascensionis quando fit Officium de qualibet Dominica minori ac de majori Dominica in Albis.}

\psalmheading{Psalmus 92 [1]}

\lettrine[lines=2]{D}{OMINUS} regnávit, decórem indútus est: * indútus est Dóminus fortitúdinem, et præcínxit se.

Etenim firmávit orbem terræ, * qui non commovébitur.

Paráta sedes tua ex tunc: * a sæculo tu es.

Elevavérunt flúmina, Dómine: * elevavérunt flúmina vocem suam.

Elevavérunt flúmina fluctus suos, * a vócibus aquárum multárum.

Mirábiles elatiónes maris: * mirábilis in altis Dóminus.

Testimónia tua credibília facta sunt nimis: * domum tuam decet sanctitúdo, Dómine, in longitúdinem diérum.

\psalmheading{Psalmus 99 [2]}

\lettrine[lines=2]{J}{UBILATE} Deo, omnis terra: * servíte Dómino in lætítia.

Introíte in conspéctu ejus, * in exsultatióne.

Scitóte quóniam Dóminus ipse est Deus: * ipse fecit nos, et non ipsi nos.

Pópulus ejus, et oves páscuæ ejus: * introíte portas ejus in confessióne, átria ejus in hymnis: confitémini illi.

Laudáte nomen ejus: quóniam suávis est Dóminus, in ætérnum misericórdia ejus, * et usque in generatiónem et generatiónem véritas ejus.

\psalmheading{Psalmus 62 [3]}

\lettrine[lines=2]{D}{EUS}, Deus meus, * ad te de luce vígilo.

Sitívit in te ánima mea, * quam multipliciter tibi caro mea.

In terra desérta, et ínvia, et inaquósa: * sic in sancto appárui tibi, ut vidérem virtútem tuam, et glóriam tuam.

Quóniam mélior est misericórdia tua super vitas: * lábia mea laudábunt te.

Sic benedícam te in vita mea: * et in nómine tuo levábo manus meas.

Sicut ádipe et pinguédine repleátur ánima mea: * et lábiis exsultatiónis laudábit os meum.

Si memor fui tui super stratum meum, in matutínis meditábor in te: * quia fuísti adjútor meus.

Et in velaménto alárum tuárum exsultábo, adhǽsit ánima mea post te: * me suscépit déxtera tua.

Ipsi vero in vanum quæsiérunt ánimam meam, introíbunt in inferióra terræ: * tradéntur in manus gládii, partes vúlpium erunt.

Rex vero lætábitur in Deo, laudabúntur omnes qui jurant in eo: * quia obstrúctum est os loquéntium iníqua.

\canticumheading{Canticum trium Puerorum [4]}{Dan. 3, 37-38 et 56}

\lettrine[lines=2]{B}{ENEDICITE}, ómnia ópera Dómini, Dómino: * laudáte et superexaltáte eum in sæcula.

Benedícite, Angeli Dómini, Dómino: * benedícite, cæli, Dómino.

Benedícite, aquæ omnes, quæ super cælos sunt, Dómino: * benedícite, omnes virtútes Dómini, Dómino.

Benedícite, sol et luna, Dómino: * benedícite, stellæ cæli, Dómino.

Benedícite, omnis imber et ros, Dómino: * benedícite, omnes spíritus Dei, Dómino.

Benedícite, ignis et æstus, Dómino: * benedícite, frigus et æstus, Dómino.

Benedícite, rores et pruína, Dómino: * benedícite, gelu et frigus, Dómino.

Benedícite, glácies et nives, Dómino: * benedícite, noctes et dies, Dómino.

Benedícite, lux et ténebræ, Dómino: * benedícite, fúlgura et nubes, Dómino.

Benedícat terra Dóminum: * laudet et superexáltet eum in sæcula.

Benedícite, montes et colles, Dómino: * benedícite, univérsa germinántia in terra, Dómino.

Benedícite, fontes, Dómino: * benedícite, mária et flúmina, Dómino.

Benedícite, cete, et ómnia, quæ movéntur in aquis, Dómino: * benedícite, omnes vólucres cæli, Dómino.

Benedícite, omnes béstiæ et pécora, Dómino: * benedícite, fílii hóminum, Dómino.

Benedícat Israël Dóminum: * laudet et superexáltet eum in sæcula.

Benedícite, sacerdótes Dómini, Dómino: * benedícite, servi Dómini, Dómino.

Benedícite, spíritus, et ánimæ justórum, Dómino: * benedícite, sancti, et húmiles corde, Dómino.

Benedícite, Ananía, Azaría, Mísaël, Dómino: * laudáte et superexaltáte eum in sæcula.

Benedicámus Patrem et Fílium cum Sancto Spíritu: * laudémus et superexaltémus eum in sæcula.

Benedíctus es, Dómine, in firmaménto cæli: * et laudábilis, et gloriósus, et superexaltátus in sæcula.

\rubric{¶ Hic non dicitur} Glória Patri, \rubric{neque} Amen.

\psalmheading{Psalmus 148 [5]}

\lettrine[lines=2]{L}{AUDATE} Dóminum de cælis: * laudáte eum in excélsis.

Laudáte eum, omnes Angeli ejus: * laudáte eum, omnes virtútes ejus.

Laudáte eum, sol et luna: * laudáte eum, omnes stellæ et lumen.

Laudáte eum, cæli cælórum: * et aquæ omnes, quæ super cælos sunt, laudent nomen Dómini.

Quia ipse dixit, et facta sunt: * ipse mandávit, et creáta sunt.

Státuit ea in ætérnum, et in sæculum sæculi: * præcéptum pósuit, et non præteríbit.

Laudáte Dóminum de terra, * dracónes, et omnes abýssi.

Ignis, grando, nix, glácies, spíritus procellárum: * quæ fáciunt verbum ejus:

Montes, et omnes colles: * ligna fructífera, et omnes cedri.

Béstiæ, et univérsa pécora: * serpéntes, et vólucres pennátæ:

Reges terræ, et omnes pópuli: * príncipes, et omnes júdices terræ.

Júvenes, et vírgines: senes cum junióribus laudent nomen Dómini: * quia exaltátum est nomen ejus solíus.

Conféssio ejus super cælum et terram: * et exaltávit cornu pópuli sui.

Hymnus ómnibus sanctis ejus: * fíliis Israël, pópulo appropinquánti sibi.

\antlabel{Ant.} Allelúja, allelúja, allelúja.

\rubric{Capitulum, Hymnus et Versus, Ant. ad Benedictus et Oratio ut in Proprio de Tempore.}

\canticumheading{Canticum Zachariæ}{Luc. 1, 68-79}

\lettrine[lines=2]{B}{ENEDICTUS} \textdagger\ Dóminus, Deus Israël, * quia visitávit, et fecit redemptiónem plebis suæ:

Et eréxit cornu salútis nobis: * in domo David, púeri sui.

Sicut locútus est per os sanctórum, * qui a século sunt, Prophetárum ejus:

Salútem ex inimícis nostris, * et de manu ómnium, qui odérunt nos:

Ad faciéndam misericórdiam cum pátribus nostris: * et memorári testaménti sui sancti.

Jusjurándum, quod jurávit ad Abraham, patrem nostrum, * datúrum se nobis:

Ut sine timóre, de manu inimicórum nostrórum liberáti, * serviámus illi.

In sanctitáte, et justítia coram ipso, * ómnibus diébus nostris.

Et tu, puer, Prophéta Altíssimi vocáberis: * præíbis enim ante fáciem Dómini paráre vias ejus:

Ad dandam sciéntiam salútis plebi ejus: * in remissiónem peccatórum eórum:

Per víscera misericórdiæ Dei nostri: * in quibus visitávit nos, óriens ex alto:

Illumináre his, qui in ténebris, et in umbra mortis sedent: * ad dirigéndos pedes nostros in viam pacis.

\rubric{Oratio conveniens.}

\rubric{Reliqua ut in Ordinario. 13}

\rubric{Psalmi cum Cantico de Dominica superius appositi dicuntur etiam, sub Antiphonis Officio diei congruentibus, in Dominica infra Octavam Ascensionis, in feria VI post Octavam Ascensionis, in Vigilia Pentecostes, in Duplicibus I vel II classis, in Festis novem Lectionum Domini, B. Mariæ Virginis, Angelorum, S. Joannis Baptistæ, S. Joseph, Apostolorum, Evangelistarum, itemque in quolibet Festo novem Lectionum proprias Antiphonas ad Laudes habenti, quando de eis persolvitur Officium.}

\hourheading{Ad Laudes II}

\textit{A Dominica I Quadragesimæ usque ad Dominicam Palmarum inclusive.}

(Pater noster \rubric{et} Ave.)

\V\ Deus, in adjutórium.

\rubric{Antiphonæ de Laudibus Dominicæ currentis, ut in Proprio de Tempore}

\psalmheading{Psalmus 50}

\lettrine[lines=2]{M}{ISERERE} mei, Deus, * secúndum magnam misericórdiam tuam.

Et secúndum multitúdinem miseratiónum tuárum, * dele iniquitátem meam.

Amplius lava me ab iniquitáte mea: * et a peccáto meo munda me.

Quóniam iniquitátem meam ego cognósco: * et peccátum meum contra me est semper.

Tibi soli peccávi, et malum coram te feci: * ut justificéris in sermónibus tuis, et vincas cum judicáris.

Ecce enim, in iniquitátibus concéptus sum: * et in peccátis concépit me mater mea.

Ecce enim, veritátem dilexísti: * incérta et occúlta sapiéntiæ tuæ manifestásti mihi.

Aspérges me hyssópo, et mundábor: * lavábis me, et super nivem dealbábor.

Audítui meo dabis gáudium et lætítiam: * et exsultábunt ossa humiliáta.

Avérte fáciem tuam a peccátis meis: * et omnes iniquitátes meas dele.

Cor mundum crea in me, Deus: * et spíritum rectum ínnova in viscéribus meis.

Ne projícias me a fácie tua: * et spíritum sanctum tuum ne áuferas a me.

Redde mihi lætítiam salutáris tui: * et spíritu principáli confírma me.

Docébo iníquos vias tuas: * et ímpii ad te converténtur.

Líbera me de sanguínibus, Deus, Deus salútis meæ: * et exsultábit lingua mea justítiam tuam.

Dómine, lábia mea apéries: * et os meum annuntiábit laudem tuam.

Quóniam si voluísses sacrifícium, dedíssem útique: * holocáustis non delectáberis.

Sacrifícium Deo spíritus contribulátus: * cor contrítum, et humiliátum, Deus, non despícies.

Benígne fac, Dómine, in bona voluntáte tua Sion: * ut ædificéntur muri Jerúsalem.

Tunc acceptábis sacrifícium justítiæ, oblatiónes, et holocáusta: * tunc impónent super altáre tuum vítulos.

\psalmheading{Psalmus 117}

\lettrine[lines=2]{C}{ONFITEMINI} Dómino quóniam bonus: * quóniam in sæculum misericórdia ejus.

Dicat nunc Israël quóniam bonus: * quóniam in sæculum misericórdia ejus.

Dicat nunc domus Aaron: * quóniam in sæculum misericórdia ejus.

Dicant nunc qui timent Dóminum: * quóniam in sæculum misericórdia ejus.

De tribulatióne invocávi Dóminum: * et exaudívit me in latitúdine Dóminus.

Dóminus mihi adjútor: * non timébo quid fáciat mihi homo.

Dóminus mihi adjútor: * et ego despíciam inimícos meos.

Bonum est confídere in Dómino, * quam confídere in hómine:

Bonum est speráre in Dómino, * quam speráre in princípibus.

Omnes gentes circuiérunt me: * et in nómine Dómini quia ultus sum in eos.

Circumdántes circumdedérunt me: * et in nómine Dómini quia ultus sum in eos.

Circumdedérunt me sicut apes, et exarsérunt sicut ignis in spinis: * et in nómine Dómini quia ultus sum in eos.

Impúlsus evérsus sum ut cáderem: * et Dóminus suscépit me.

Fortitúdo mea, et laus mea Dóminus: * et factus est mihi in salútem.

Vox exsultatiónis, et salútis * in tabernáculis justórum.

Déxtera Dómini fecit virtútem: déxtera Dómini exaltávit me, * déxtera Dómini fecit virtútem.

Non móriar, sed vivam: * et narrábo ópera Dómini.

Castígans castigávit me Dóminus: * et morti non trádidit me.

Aperíte mihi portas justítiæ, ingréssus in eas confitébor Dómino: * hæc porta Dómini, justi intrábunt in eam.

Confitébor tibi quóniam exaudísti me: * et factus es mihi in salútem.

Lápidem, quem reprobavérunt ædificántes: * hic factus est in caput ánguli.

A Dómino factum est istud: * et est mirábile in óculis nostris.

Hæc est dies, quam fecit Dóminus: * exsultémus, et lætémur in ea.

O Dómine, salvum me fac, o Dómine, bene prosperáre: * benedíctus qui venit in nómine Dómini.

Benedíximus vobis de domo Dómini: * Deus Dóminus, et illúxit nobis.

Constitúite diem solémnem in condénsis, * usque ad cornu altáris.

Deus meus es tu, et confitébor tibi: * Deus meus es tu, et exaltábo te.

Confitébor tibi quóniam exaudísti me: * et factus es mihi in salútem.

Confitémini Dómino quóniam bonus: * quóniam in sæculum misericórdia ejus.

\psalmheading{Psalmus 62}

\lettrine[lines=2]{D}{EUS}, Deus meus, * ad te de luce vígilo.

Sitívit in te ánima mea, * quam multipliciter tibi caro mea.

In terra desérta, et ínvia, et inaquósa: * sic in sancto appárui tibi, ut vidérem virtútem tuam, et glóriam tuam.

Quóniam mélior est misericórdia tua super vitas: * lábia mea laudábunt te.

Sic benedícam te in vita mea: * et in nómine tuo levábo manus meas.

Sicut ádipe et pinguédine repleátur ánima mea: * et lábiis exsultatiónis laudábit os meum.

Si memor fui tui super stratum meum, in matutínis meditábor in te: * quia fuísti adjútor meus.

Et in velaménto alárum tuárum exsultábo, adhǽsit ánima mea post te: * me suscépit déxtera tua.

Ipsi vero in vanum quæsiérunt ánimam meam, introíbunt in inferióra terræ: * tradéntur in manus gládii, partes vúlpium erunt.

Rex vero lætábitur in Deo, laudabúntur omnes qui jurant in eo: * quia obstrúctum est os loquéntium iníqua.

\canticumheading{Canticum trium Puerorum}{Dan. 3, 53-57}

\lettrine[lines=2]{B}{ENEDICTUS} es, Dómine, Deus patrum nostrórum: * et laudábilis, et gloriósus, et superexaltátus in sæcula.

Et benedíctum nomen glóriæ tuæ sanctum: * et laudábile, et superexaltátum in ómnibus sæculis.

Benedíctus es in templo sancto glóriæ tuæ: * et superlaudábilis, et supergloriósus in sæcula.

Benedíctus es in throno regni tui: * et superlaudábilis, et superexaltátus in sæcula.

Benedíctus es, qui intuéris abýssos, et sedes super Chérubim: * et laudábilis, et superexaltátus in sæcula.

Benedíctus es in firmaménto cæli: * et laudábilis, et gloriósus in sæcula.

Benedícite, ómnia ópera Dómini, Dómino: * laudáte, et superexaltáte eum in sæcula.

Glória Patri.

\psalmheading{Psalmus 148}

\lettrine[lines=2]{L}{AUDATE} Dóminum de cælis: * laudáte eum in excélsis.

Laudáte eum, omnes Angeli ejus: * laudáte eum, omnes virtútes ejus.

Laudáte eum, sol et luna: * laudáte eum, omnes stellæ et lumen.

Laudáte eum, cæli cælórum: * et aquæ omnes, quæ super cælos sunt, laudent nomen Dómini.

Quia ipse dixit, et facta sunt: * ipse mandávit, et creáta sunt.

Státuit ea in ætérnum, et in sæculum sæculi: * præcéptum pósuit, et non præteríbit.

Laudáte Dóminum de terra, * dracónes, et omnes abýssi.

Ignis, grando, nix, glácies, spíritus procellárum: * quæ fáciunt verbum ejus:

Montes, et omnes colles: * ligna fructífera, et omnes cedri.

Béstiæ, et univérsa pécora: * serpéntes, et vólucres pennátæ:

Reges terræ, et omnes pópuli: * príncipes, et omnes júdices terræ.

Júvenes, et vírgines: senes cum junióribus laudent nomen Dómini: * quia exaltátum est nomen ejus solíus.

Conféssio ejus super cælum, et terram: * et exaltávit cornu pópuli sui.

Hymnus ómnibus sanctis ejus: * fíliis Israël, pópulo appropinquánti sibi.

\rubric{Capitulum, Hymnus, Versus, Antiphona ad Benedictus et Oratio ut in Proprio de Tempore.}

\rubric{Reliqua ut in Ordinario. 13}

% ===== END OF BATCH 3 (Dominica, Ad Laudes I/II; bare trailing "13"
% after "Reliqua ut in Ordinario." at verna.md lines 2411 and 2601
% preserved literally, not resolved -- no page-marker \label exists
% in this file's Ordinarium to resolve it to, and it is not a
% bracketed [REF] tag; same treatment as the bare "[3]" markers after
% "Ps. 94 Veníte" in the Ordinarium and Matutinum), continues with
% Ad Primam =====

% ===================================================================
% BATCH 4: Psalterium, Dominica -- Ad Primam, Ad Tertiam, Ad Sextam,
% Ad Nonam (verna.md lines 2603-3424)
% ===================================================================

\hourheading{Ad Primam}

Pater, Ave et Credo.

\V\ Deus, in adjutórium.

\hymnusheading{Hymnus}

\lettrine[lines=2]{J}{AM} lucis orto sídere,\\
Deum precémur súpplices,\\
Ut in diúrnis áctibus\\
Nos servet a nocéntibus.

Linguam refrénans témperet,\\
Ne litis horror ínsonet:\\
Visum fovéndo cóntegat,\\
Ne vanitátes háuriat.

Sint pura cordis íntima,\\
Absístat et vecórdia:\\
Carnis terat supérbiam\\
Potus cibíque párcitas.

Ut, cum dies abscésserit\\
Noctémque sors redúxerit,\\
Mundi per abstinéntiam\\
Ipsi canámus glóriam.

Deo Patri sit glória,\\
Ejúsque soli Fílio,\\
Cum Spíritu Paráclito,\\
Nunc et per omne sæculum. Amen.

\antlabel{Ant.} Allelúja.

\rubric{A Dominica I Quadragesimæ usque ad Dominicam Palmarum inclusive: Ant. ut in Proprio de Tempore; atque loco sequentis Ps. 117 Confitémini, dicuntur sequ. Ps. 92 et 99, qui in Laudibus locum non habuerunt.}

\psalmheading{Psalmus 92}

\lettrine[lines=2]{D}{OMINUS} regnávit, decórem indútus est: * indútus est Dóminus fortitúdinem, et præcínxit se.

Etenim firmávit orbem terræ, * qui non commovébitur.

Paráta sedes tua ex tunc: * a sæculo tu es.

Elevavérunt flúmina, Dómine: * elevavérunt flúmina vocem suam.

Elevavérunt flúmina fluctus suos, * a vócibus aquárum multárum.

Mirábiles elatiónes maris: * mirábilis in altis Dóminus.

Testimónia tua credibília facta sunt nimis: * domum tuam decet sanctitúdo, Dómine, in longitúdinem diérum.

\psalmheading{Psalmus 99}

\lettrine[lines=2]{J}{UBILATE} Deo, omnis terra: * servíte Dómino in lætítia.

Introíte in conspéctu ejus, * in exsultatióne.

Scitóte quóniam Dóminus ipse est Deus: * ipse fecit nos, et non ipsi nos.

Pópulus ejus, et oves páscuæ ejus: * introíte portas ejus in confessióne, átria ejus in hymnis: confitémini illi.

Laudáte nomen ejus: quóniam suávis est Dóminus, in ætérnum misericórdia ejus, * et usque in generatiónem et generatiónem véritas ejus.

\variantheading{Tempore Paschali}

\antlabel{Ant.} Allelúja.

\rubric{¶ Antiphona cum Psalmis Tempore Paschali ad Horas minores assignata, adhibetur semper, extra Octavam Ascensionis, quando fit Officium de qualibet Dominica minori ac de majori Dominica in Albis.}

\psalmheading{Psalmus 117}

\lettrine[lines=2]{C}{ONFITEMINI} Dómino quóniam bonus: * quóniam in sæculum misericórdia ejus.

Dicat nunc Israël quóniam bonus: * quóniam in sæculum misericórdia ejus.

Dicat nunc domus Aaron: * quóniam in sæculum misericórdia ejus.

Dicant nunc qui timent Dóminum: * quóniam in sæculum misericórdia ejus.

De tribulatióne invocávi Dóminum: * et exaudívit me in latitúdine Dóminus.

Dóminus mihi adjútor: * non timébo quid fáciat mihi homo.

Dóminus mihi adjútor: * et ego despíciam inimícos meos.

Bonum est confídere in Dómino, * quam confídere in hómine:

Bonum est speráre in Dómino, * quam speráre in princípibus.

Omnes gentes circuiérunt me: * et in nómine Dómini quia ultus sum in eos.

Circumdántes circumdedérunt me: * et in nómine Dómini quia ultus sum in eos.

Circumdedérunt me sicut apes, et exarsérunt sicut ignis in spinis: * et in nómine Dómini quia ultus sum in eos.

Impúlsus evérsus sum ut cáderem: * et Dóminus suscépit me.

Fortitúdo mea, et laus mea Dóminus: * et factus est mihi in salútem.

Vox exsultatiónis, et salútis * in tabernáculis justórum.

Déxtera Dómini fecit virtútem: déxtera Dómini exaltávit me, * déxtera Dómini fecit virtútem.

Non móriar, sed vivam: * et narrábo ópera Dómini.

Castígans castigávit me Dóminus: * et morti non trádidit me.

Aperíte mihi portas justítiæ, ingréssus in eas confitébor Dómino: * hæc porta Dómini, justi intrábunt in eam.

Confitébor tibi quóniam exaudísti me: * et factus es mihi in salútem.

Lápidem, quem reprobavérunt ædificántes: * hic factus est in caput ánguli.

A Dómino factum est istud: * et est mirábile in óculis nostris.

Hæc est dies, quam fecit Dóminus: * exsultémus, et lætémur in ea.

O Dómine, salvum me fac, o Dómine, bene prosperáre: * benedíctus qui venit in nómine Dómini.

Benedíximus vobis de domo Dómini: * Deus Dóminus, et illúxit nobis.

Constitúite diem solémnem in condénsis, * usque ad cornu altáris.

Deus meus es tu, et confitébor tibi: * Deus meus es tu, et exaltábo te.

Confitébor tibi quóniam exaudísti me: * et factus es mihi in salútem.

Confitémini Dómino quóniam bonus: * quóniam in sæculum misericórdia ejus.

\rubric{¶ In ultimo triduo Majoris Hebdomadæ, in Dominica infra Octavam Ascensionis, in Feria VI post Octavam Ascensionis, in Vigilia Pentecostes, in Duplicibus I vel II classis, in Festis novem Lectionum Domini, beatæ Mariæ Virginis, Angelorum, sancti Joannis Baptistæ, sancti Joseph, Apostolorum, Evangelistarum, et per Octavas privilegiatas; si de eis fiat Officium, ad Primam, Tertiam, Sextam et Nonam dicuntur Psalmi de Dominica, sub Antiphonis Officio diei congruentibus, nisi aliter suo loco notetur; sed ad Primam, loco Psalmi 117 Confitémini, in ipsa etiam die Dominica, semper dicitur}

\psalmheading{Psalmus 53}

\lettrine[lines=2]{D}{EUS}, in nómine tuo salvum me fac: * et in virtúte tua júdica me.

Deus, exáudi oratiónem meam: * áuribus pércipe verba oris mei.

Quóniam aliéni insurrexérunt advérsum me, et fortes quæsiérunt ánimam meam: * et non proposuérunt Deum ante conspéctum suum.

Ecce enim, Deus ádjuvat me: * et Dóminus suscéptor est ánimæ meæ.

Avérte mala inimícis meis: * et in veritáte tua dispérde illos.

Voluntárie sacrificábo tibi, * et confitébor nómini tuo, Dómine: quóniam bonum est:

Quóniam ex omni tribulatióne eripuísti me: * et super inimícos meos despéxit óculus meus.

\psalmheading{Psalmus 118, i}

\lettrine[lines=2]{B}{EATI} immaculáti in via: * qui ámbulant in lege Dómini.

Beáti, qui scrutántur testimónia ejus: * in toto corde exquírunt eum.

Non enim qui operántur iniquitátem: * in viis ejus ambulavérunt.

Tu mandásti * mandáta tua custodíri nimis.

Utinam dirigántur viæ meæ, * ad custodiéndas justificatiónes tuas!

Tunc non confúndar, * cum perspéxero in ómnibus mandátis tuis.

Confitébor tibi in directióne cordis: * in eo quod dídici judícia justítiæ tuæ.

Justificatiónes tuas custódiam: * non me derelínquas usquequáque.

In quo córrigit adolescéntior viam suam? * in custodiéndo sermónes tuos.

In toto corde meo exquisívi te: * ne repéllas me a mandátis tuis.

In corde meo abscóndi elóquia tua: * ut non peccem tibi.

Benedíctus es, Dómine: * doce me justificatiónes tuas.

In lábiis meis, * pronuntiávi ómnia judícia oris tui.

In via testimoniórum tuórum delectátus sum, * sicut in ómnibus divítiis.

In mandátis tuis exercébor: * et considerábo vias tuas.

In justificatiónibus tuis meditábor: * non oblivíscar sermónes tuos.

\psalmheading{Psalmus 118, ii}

\lettrine[lines=2]{R}{ETRIBUE} servo tuo, vivífica me: * et custódiam sermónes tuos.

Revéla óculos meos: * et considerábo mirabília de lege tua.

Íncola ego sum in terra: * non abscóndas a me mandáta tua.

Concupívit ánima mea desideráre justificatiónes tuas, * in omni témpore.

Increpásti supérbos: * maledícti qui declínant a mandátis tuis.

Aufer a me oppróbrium, et contémptum: * quia testimónia tua exquisívi.

Etenim sedérunt príncipes, et advérsum me loquebántur: * servus autem tuus exercebátur in justificatiónibus tuis.

Nam et testimónia tua meditátio mea est: * et consílium meum justificatiónes tuæ.

Adhǽsit paviménto ánima mea: * vivífica me secúndum verbum tuum.

Vias meas enuntiávi, et exaudísti me: * doce me justificatiónes tuas.

Viam justificatiónum tuárum ínstrue me: * et exercébor in mirabílibus tuis.

Dormitávit ánima mea præ tædio: * confírma me in verbis tuis.

Viam iniquitátis ámove a me: * et de lege tua miserére mei.

Viam veritátis elégi: * judícia tua non sum oblítus.

Adhǽsi testimóniis tuis, Dómine: * noli me confúndere.

Viam mandatórum tuórum cucúrri, * cum dilatásti cor meum.

\rubric{Extra Tempus Paschale: Antiphona ut in Proprio de Tempore.}

\variantheading{Tempore Paschali}

\antlabel{Ant.} Allelúja, allelúja, allelúja.

\capitulumlabel{Capitulum}{1 Tim. 1, 17}

\lettrine[lines=2]{R}{EGI} sæculórum immortáli et invisíbili, soli Deo honor et glória in sæcula sæculórum. Amen.

\R\ br. Christe, Fili Dei vivi, * Miserére nobis. Christe. \V\ \textit{Qui sedes ad déxteram Patris.} Miserére nobis. Glória Patri. Christe. \V\ Exsúrge, Christe, ádjuva nos. \R\ Et líbera nos propter nomen tuum.

\rubric{In Dominicis Passionis et Palmarum omittitur} Glória Patri.

\variantheading{Tempore Paschali}

\R\ br. Christe, Fili Dei vivi, miserére nobis, * Allelúja, allelúja. Christe. \V\ \textit{Qui surrexísti a mórtuis.} Allelúja, allelúja. Glória Patri. Christe. \V\ Exsúrge, Christe, ádjuva nos, allelúja. \R\ Et líbera nos propter nomen tuum, allelúja.

\variantheading{Preces (si dicendæ sint)}

KYRIE, eléison. Christe, eléison. Kýrie, eléison.

Pater noster \rubric{secreto usque ad} \V\ Et ne nos indúcas in tentatiónem. \R\ Sed líbera nos a malo.

Credo in Deum \rubric{secreto usque ad} \V\ Carnis resurrectiónem. \R\ Vitam ætérnam. Amen. \V\ Et ego ad te, Dómine, clamávi. \R\ Et mane orátio mea prævéniet te. \V\ Repleátur os meum laude. \R\ Ut cantem glóriam tuam, tota die magnitúdinem tuam. \V\ Dómine, avérte fáciem tuam a peccátis meis. \R\ Et omnes iniquitátes meas dele. \V\ Cor mundum crea in me, Deus. \R\ Et spíritum rectum ínnova in viscéribus meis. \V\ Ne projícias me a fácie tua. \R\ Et spíritum sanctum tuum ne áuferas a me. \V\ Redde mihi lætítiam salutáris tui. \R\ Et spíritu principáli confírma me. \V\ Adjutórium nostrum in nómine Dómini. \R\ Qui fecit cælum et terram.

Confíteor. Misereátur. Indulgéntiam.

\V\ Dignáre, Dómine, die isto. \R\ Sine peccáto nos custodíre. \V\ Miserére nostri, Dómine. \R\ Miserére nostri. \V\ Fiat misericórdia tua, Dómine, super nos. \R\ Quemádmodum sperávimus in te. \V\ Dómine, exáudi oratiónem meam. \R\ Et clamor meus ad te véniat.

\V\ Dóminus vobíscum.

Orémus. \oratiolabel{Oratio}

\lettrine[lines=2]{D}{OMINE}, Deus omnípotens, qui ad princípium hujus diéi nos perveníre fecísti: tua nos hódie salva virtúte; ut in hac die ad nullum declinémus peccátum, sed semper ad tuam justítiam faciéndam nostra procédant elóquia, dirigántur cogitatiónes et ópera. Per Dóminum. \R\ Amen.

\V\ Dóminus vobíscum.

\V\ Benedicámus Dómino.

\rubric{Deinde in Choro legitur Martyrologium, quod laudabiliter fit etiam extra Chorum. Postea Hebdomadarius dicit:}

\V\ Pretiósa in conspéctu Dómini. \R\ Mors Sanctórum ejus.

\rubric{Et absolute, sine} Orémus:

\oratiolabel{Oratio}

\lettrine[lines=2]{S}{ANCTA} María, et omnes Sancti intercédant pro nobis ad Dóminum, ut nos mereámur ab eo adjuvári et salvári, qui vivit et regnat in sæcula sæculórum. \R\ Amen.

\V\ Deus, in adjutórium meum inténde. \R\ Dómine, ad adjuvándum me festína. \rubric{Et dicitur ter; ultimo additur:} Glória Patri. Sicut erat.

Kýrie, eléison. Christe, eléison. Kýrie, eléison.

Pater noster \rubric{secreto usque ad} \V\ Et ne nos indúcas in tentatiónem. \R\ Sed líbera nos a malo. \V\ Réspice in servos tuos, Dómine, et in ópera tua, et dírige fílios eórum. \R\ Et sit splendor Dómini, Dei nostri, super nos, et ópera mánuum nostrárum dírige super nos, et opus mánuum nostrárum dírige. \V\ Glória Patri. \R\ Sicut erat.

Orémus. \oratiolabel{Oratio}

\lettrine[lines=2]{D}{IRIGERE} et sanctificáre, régere et gubernáre dignáre, Dómine Deus, Rex cæli et terræ, hódie corda et córpora nostra, sensus, sermónes et actus nostros in lege tua, et in opéribus mandatórum tuórum: ut hic, et in ætérnum, te auxiliánte, salvi et líberi esse mereámur, Salvátor mundi: Qui vivis et regnas in sæcula sæculórum. \R\ Amen.

\V\ Jube, domne, benedícere. \benedlabel{Bened.} Dies et actus nostros in sua pace dispónat Dóminus omnípotens. \R\ Amen.

\variantheading{Tempore Quadragesimæ}

\lettrine[lines=2]{Q}{UÆRITE} Dóminum, dum inveníri potest: invocáte eum, dum prope est. Tu autem.

\variantheading{Tempore Passionis}

\lettrine[lines=2]{F}{ACIEM} meam non avérti ab increpántibus, et conspuéntibus in me. Dóminus Deus auxiliátor meus, et ídeo non sum confúsus. Tu autem.

\variantheading{Tempore Paschali}

\lettrine[lines=2]{S}{I} consurrexístis cum Christo, quæ sursum sunt quærite, ubi Christus est in déxtera Dei sedens: quæ sursum sunt sápite, non quæ super terram.

\sectionheading{E Proprio de Tempore}

\rubric{In Paschate et per Octavam:} Si consurrexístis, ut supra Tempore Paschali.

\feastheading{In Ascensione Domini et per Octavam}

\lettrine[lines=2]{V}{IRI} Galilǽi, quid statis aspiciéntes in cælum? Hic Jesus, qui assúmptus est a vobis in cælum, sic véniet, quemádmodum vidístis eum eúntem in cælum.

\feastheading{Dom. infra Oct. Ascens., Feria VI post Octavam et in Vigilia Pentecostes}

\lettrine[lines=2]{S}{I} quis lóquitur, quasi sermónes Dei; si quis minístrat, tamquam ex virtúte, quam admínistrat Deus: ut in ómnibus honorificétur Deus per Jesum Christum, Dóminum nostrum.

\feastheading{In Dominica Pentecostes et per Octavam}

\lettrine[lines=2]{J}{UDÆI} quoque et Prosélyti, Cretes et Arabes: audívimus eos loquéntes nostris linguis magnália Dei.

\sectionheading{E Proprio Sanctorum}

\feastheading{In Apparitione B. M. V.}

\lettrine[lines=2]{I}{N} me grátia omnis viæ et veritátis, in me omnis spes vitæ et virtútis. Transíte ad me, omnes qui concupíscitis me, et a generatiónibus meis implémini.

\feastheading{S. Joseph 19 Mart.}

\lettrine[lines=2]{P}{ROFUGUM} justum dedúxit Sapiéntia per vias rectas, et osténdit illi regnum Dei, et dedit illi sciéntiam sanctórum: honestávit illum in labóribus, et complévit labóres illíus.

\feastheading{S. Gabrielis Archangeli}

\lettrine[lines=2]{S}{CITO} ergo et animadvérte: Ab éxitu sermónis, ut íterum ædificétur Jerúsalem, usque ad Christum ducem, hebdómades septem et hebdómades sexagínta duæ erunt.

\feastheading{In Annuntiatione B. M. V.}

\lettrine[lines=2]{E}{GREDIETUR} virga de radíce Jesse, et flos de radíce ejus ascéndet. Et requiéscet super eum Spíritus Dómini.

\feastheading{Septem Dolorum B. M. V.}

\lettrine[lines=2]{G}{ENERATIONEM} ejus quis enarrábit? quia abscíssus est de terra vivéntium: propter scelus pópuli mei percússi eum.

\feastheading{In Solemnitate S. Joseph}

\lettrine[lines=2]{F}{ILIUS} accréscens Joseph, fílius accréscens, et decórus aspéctu: fíliæ discurrérunt super murum.

\feastheading{In Inventione S. Crucis}

\lettrine[lines=2]{H}{UMILIAVIT} semetípsum factus obédiens usque ad mortem, mortem autem Crucis. Propter quod et Deus exaltávit illum, et donávit illi nomen, quod est super omne nomen.

\feastheading{In Apparit. S. Michaelis}

\lettrine[lines=2]{F}{ACTUM} est prœlium magnum in cælo: Míchaël et Angeli ejus prœliabántur cum dracóne, et draco pugnábat et ángeli ejus; et non prævaluérunt, neque locus invéntus est eórum ámplius in cælo.

\sectionheading{E Communi Sanctorum}

\feastheading{Pro Ap. et Ev., (extra T. P.)}

\lettrine[lines=2]{I}{BANT} Apóstoli gaudéntes a conspéctu concílii, quóniam digni habíti sunt pro nómine Jesu contuméliam pati.

\feastheading{Pro Ap. et Ev., uno et plur. Mm. T. P.}

\lettrine[lines=2]{S}{CIMUS} quóniam diligéntibus Deum ómnia cooperántur in bonum, iis qui secúndum propósitum vocáti sunt sancti.

\feastheading{Pro Martyre (extra T. P.)}

\lettrine[lines=2]{J}{USTUS} cor suum trádidit ad vigilándum dilúculo ad Dóminum, qui fecit illum, et in conspéctu Altíssimi deprecábitur.

\feastheading{Pro plur. Mm. (extra T. P.)}

\lettrine[lines=2]{F}{ULGEBUNT} justi, et tamquam scintíllæ in arundinéto discúrrent. Judicábunt natiónes, et dominabúntur pópulis: et regnábit Dóminus illórum in perpétuum.

\feastheading{Pro Conf. Pont.}

\lettrine[lines=2]{F}{UNGI} sacerdótio, et habére laudem in nómine ipsíus, et offérre illi incénsum dignum in odórem suavitátis.

\feastheading{Pro Conf. non Pont.}

\lettrine[lines=2]{J}{USTUM} dedúxit Dóminus per vias rectas, et osténdit illi regnum Dei, et dedit illi sciéntiam sanctórum; honestávit illum in labóribus, et complévit labóres illíus.

\feastheading{Pro Virgine Martyre}

\lettrine[lines=2]{D}{OMINE}, Deus meus, exaltásti super terram habitatiónem meam, et pro morte defluénte deprecáta sum. Invocávi Dóminum, Patrem Dómini mei, ut non derelínquat me in die tribulatiónis meæ, et in témpore superbórum sine adjutório.

\feastheading{Pro Virgine non Martyre}

\lettrine[lines=2]{O}{} \textsc{quam} pulchra est casta generátio cum claritáte! immortális est enim memória illíus: quóniam apud Deum nota est, et apud hómines.

\feastheading{Pro Martyre tantum}

\lettrine[lines=2]{L}{AUDABIT} usque ad mortem ánima mea Dóminum, quóniam éruis sustinéntes te, et líberas eos de manu angústiæ, Dómine, Deus noster.

\feastheading{Pro nec Virg. nec Mart.}

\lettrine[lines=2]{M}{ULTAE} fíliæ congregavérunt divítias: tu supergréssa es univérsas. Fallax grátia, et vana est pulchritúdo: múlier timens Dóminum, ipsa laudábitur.

\feastheading{Pro Dedicatione Ecclesiæ}

\lettrine[lines=2]{E}{T} abstérget Deus omnem lácrimam ab óculis eórum: et mors ultra non erit, neque luctus, neque clamor, neque dolor erit ultra, quia prima abiérunt. Et dixit qui sedébat in throno: Ecce, nova fácio ómnia.

\feastheading{In Festis B. Mariæ V.}

\lettrine[lines=2]{I}{N} platéis sicut cinnamómum et bálsamum aromatízans odórem dedi: quasi myrrha elécta dedi suavitátem odóris.

\V\ Adjutórium nostrum in nómine Dómini. \R\ Qui fecit cælum et terram.

\V\ Benedícite. \R\ Deus.

\benedictiolabel{Benedictio} Dóminus nos benedícat, et ab omni malo deféndat, et ad vitam perdúcat ætérnam. Et fidélium ánimæ per misericórdiam Dei requiéscant in pace. \R\ Amen.

\rubric{Deinde dicitur tantum} Pater noster \rubric{secreto, nisi sequatur alia Hora.}

\hourheading{Ad Tertiam}

Pater et Ave.

\V\ Deus, in adjutórium.

\hymnusheading{Hymnus}

\lettrine[lines=2]{N}{UNC}, Sancte, nobis, Spíritus,\\
Unum Patri cum Fílio,\\
Dignáre promptus íngeri\\
Nostro refúsus péctori.

Os, lingua, mens, sensus, vigor\\
Confessiónem pérsonent,\\
Flamméscat igne cáritas,\\
Accéndat ardor próximos.

* Præsta, Pater piíssime,\\
Patríque compar Unice,\\
Cum Spíritu Paráclito\\
Regnans per omne sæculum. Amen.

\rubric{¶ In Festo et per Octavam Pentecostes, loco præcedentis, dicitur Hymnus} Veni Creátor, \rubric{ut in Proprio de Tempore.}

\rubric{Extra Tempus Paschale: Antiphona ut in Proprio de Tempore.}

\variantheading{Tempore Paschali}

\antlabel{Ant.} Allelúja.

\psalmheading{Psalmus 118, iii}

\lettrine[lines=2]{L}{EGEM} pone mihi, Dómine, viam justificatiónum tuárum: * et exquíram eam semper.

Da mihi intelléctum, et scrutábor legem tuam: * et custódiam illam in toto corde meo.

Deduc me in sémitam mandatórum tuórum: * quia ipsam vólui.

Inclína cor meum in testimónia tua: * et non in avarítiam.

Avérte óculos meos ne vídeant vanitátem: * in via tua vivífica me.

Státue servo tuo elóquium tuum, * in timóre tuo.

Amputa oppróbrium meum quod suspicátus sum: * quia judícia tua jucúnda.

Ecce, concupívi mandáta tua: * in æquitáte tua vivífica me.

Et véniat super me misericórdia tua, Dómine: * salutáre tuum secúndum elóquium tuum.

Et respondébo exprobrántibus mihi verbum: * quia sperávi in sermónibus tuis.

Et ne áuferas de ore meo verbum veritátis usquequáque: * quia in judíciis tuis supersperávi.

Et custódiam legem tuam semper: * in sæculum et in sæculum sæculi.

Et ambulábam in latitúdine: * quia mandáta tua exquisívi.

Et loquébar in testimóniis tuis in conspéctu regum: * et non confundébar.

Et meditábar in mandátis tuis, * quæ diléxi.

Et levávi manus meas ad mandáta tua, quæ diléxi: * et exercébar in justificatiónibus tuis.

\psalmheading{Psalmus 118, iv}

\lettrine[lines=2]{M}{EMOR} esto verbi tui servo tuo, * in quo mihi spem dedísti.

Hæc me consoláta est in humilitáte mea: * quia elóquium tuum vivificábit me.

Supérbi iníque agébant usquequáque: * a lege autem tua non declinávi.

Memor fui judiciórum tuórum a sæculo, Dómine: * et consolátus sum.

Deféctio ténuit me, * pro peccatóribus derelinquéntibus legem tuam.

Cantábiles mihi erant justificatiónes tuæ, * in loco peregrinatiónis meæ.

Memor fui nocte nóminis tui, Dómine: * et custodívi legem tuam.

Hæc facta est mihi: * quia justificatiónes tuas exquisívi.

Pórtio mea, Dómine, * dixi custodíre legem tuam.

Deprecátus sum fáciem tuam in toto corde meo: * miserére mei secúndum elóquium tuum.

Cogitávi vias meas: * et convérti pedes meos in testimónia tua.

Parátus sum, et non sum turbátus: * ut custódiam mandáta tua.

Funes peccatórum circumpléxi sunt me: * et legem tuam non sum oblítus.

Média nocte surgébam ad confiténdum tibi, * super judícia justificatiónis tuæ.

Párticeps ego sum ómnium timéntium te: * et custodiéntium mandáta tua.

Misericórdia tua, Dómine, plena est terra: * justificatiónes tuas doce me.

\psalmheading{Psalmus 118, v}

\lettrine[lines=2]{B}{ONITATEM} fecísti cum servo tuo, Dómine, * secúndum verbum tuum.

Bonitátem, et disciplínam, et sciéntiam doce me: * quia mandátis tuis crédidi.

Priúsquam humiliárer ego delíqui: * proptérea elóquium tuum custodívi.

Bonus es tu: * et in bonitáte tua doce me justificatiónes tuas.

Multiplicáta est super me iníquitas superbórum: * ego autem in toto corde meo scrutábor mandáta tua.

Coagulátum est sicut lac cor eórum: * ego vero legem tuam meditátus sum.

Bonum mihi quia humiliásti me: * ut discam justificatiónes tuas.

Bonum mihi lex oris tui, * super míllia auri et argénti.

Manus tuæ fecérunt me, et plasmavérunt me: * da mihi intelléctum, et discam mandáta tua.

Qui timent te vidébunt me, et lætabúntur: * quia in verba tua supersperávi.

Cognóvi, Dómine, quia æquitas judícia tua: * et in veritáte tua humiliásti me.

Fiat misericórdia tua ut consolétur me, * secúndum elóquium tuum servo tuo.

Véniant mihi miseratiónes tuæ, et vivam: * quia lex tua meditátio mea est.

Confundántur supérbi, quia injúste iniquitátem fecérunt in me: * ego autem exercébor in mandátis tuis.

Convertántur mihi timéntes te: * et qui novérunt testimónia tua.

Fiat cor meum immaculátum in justificatiónibus tuis, * ut non confúndar.

\rubric{Extra Tempus Paschale: Antiphona ut in Proprio de Tempore.}

\variantheading{Tempore Paschali}

\antlabel{Ant.} Allelúja, allelúja, allelúja.

\rubric{Reliqua ut in Proprio de Tempore et in Ordinario. 23}

\hourheading{Ad Sextam}

Pater noster et Ave.

\V\ Deus, in adjutórium.

\hymnusheading{Hymnus}

\lettrine[lines=2]{R}{ECTOR} potens, verax Deus,\\
Qui témperas rerum vices,\\
Splendóre mane illúminas,\\
Et ígnibus merídiem:

Exstíngue flammas lítium,\\
Aufer calórem nóxium,\\
Confer salútem córporum,\\
Verámque pacem córdium.

* Præsta, Pater piíssime,\\
Patríque compar Unice,\\
Cum Spíritu Paráclito\\
Regnans per omne sæculum. Amen.

\rubric{Extra Tempus Paschale: Antiphona ut in Proprio de Tempore.}

\variantheading{Tempore Paschali}

\antlabel{Ant.} Allelúja.

\psalmheading{Psalmus 118, vi}

\lettrine[lines=2]{D}{EFECIT} in salutáre tuum ánima mea: * et in verbum tuum supersperávi.

Defecérunt óculi mei in elóquium tuum, * dicéntes: Quando consoláberis me?

Quia factus sum sicut uter in pruína: * justificatiónes tuas non sum oblítus.

Quot sunt dies servi tui? * quando fácies de persequéntibus me judícium?

Narravérunt mihi iníqui fabulatiónes: * sed non ut lex tua.

Omnia mandáta tua véritas: * iníque persecúti sunt me, ádjuva me.

Paulo minus consummavérunt me in terra: * ego autem non derelíqui mandáta tua.

Secúndum misericórdiam tuam vivífica me: * et custódiam testimónia oris tui.

In ætérnum, Dómine, * verbum tuum pérmanet in cælo.

In generatiónem et generatiónem véritas tua: * fundásti terram, et pérmanet.

Ordinatióne tua persevérat dies: * quóniam ómnia sérviunt tibi.

Nisi quod lex tua meditátio mea est: * tunc forte periíssem in humilitáte mea.

In ætérnum non oblivíscar justificatiónes tuas: * quia in ipsis vivificásti me.

Tuus sum ego, salvum me fac: * quóniam justificatiónes tuas exquisívi.

Me exspectavérunt peccatóres ut pérderent me: * testimónia tua intelléxi.

Omnis consummatiónis vidi finem: * latum mandátum tuum nimis.

\psalmheading{Psalmus 118, vii}

\lettrine[lines=2]{Q}{UOMODO} diléxi legem tuam, Dómine? * tota die meditátio mea est.

Super inimícos meos prudéntem me fecísti mandáto tuo: * quia in ætérnum mihi est.

Super omnes docéntes me intelléxi: * quia testimónia tua meditátio mea est.

Super senes intelléxi: * quia mandáta tua quæsívi.

Ab omni via mala prohíbui pedes meos: * ut custódiam verba tua.

A judíciis tuis non declinávi: * quia tu legem posuísti mihi.

Quam dúlcia fáucibus meis elóquia tua, * super mel ori meo!

A mandátis tuis intelléxi: * proptérea odívi omnem viam iniquitátis.

Lucérna pédibus meis verbum tuum, * et lumen sémitis meis.

Jurávi, et státui * custodíre judícia justítiæ tuæ.

Humiliátus sum usquequáque, Dómine: * vivífica me secúndum verbum tuum.

Voluntária oris mei beneplácita fac, Dómine: * et judícia tua doce me.

Anima mea in mánibus meis semper: * et legem tuam non sum oblítus.

Posuérunt peccatóres láqueum mihi: * et de mandátis tuis non errávi.

Hereditáte acquisívi testimónia tua in ætérnum: * quia exsultátio cordis mei sunt.

Inclinávi cor meum ad faciéndas justificatiónes tuas in ætérnum, * propter retributiónem.

\psalmheading{Psalmus 118, viii}

\lettrine[lines=2]{I}{NIQUOS} ódio hábui: * et legem tuam diléxi.

Adjútor et suscéptor meus es tu: * et in verbum tuum supersperávi.

Declináte a me, malígni: * et scrutábor mandáta Dei mei.

Súscipe me secúndum elóquium tuum, et vivam: * et non confúndas me ab exspectatióne mea.

Adjuva me, et salvus ero: * et meditábor in justificatiónibus tuis semper.

Sprevísti omnes discedéntes a judíciis tuis: * quia injústa cogitátio eórum.

Prævaricántes reputávi omnes peccatóres terræ: * ídeo diléxi testimónia tua.

Confíge timóre tuo carnes meas: * a judíciis enim tuis tímui.

Feci judícium et justítiam: * non tradas me calumniántibus me.

Súscipe servum tuum in bonum: * non calumniéntur me supérbi.

Óculi mei defecérunt in salutáre tuum: * et in elóquium justítiæ tuæ.

Fac cum servo tuo secúndum misericórdiam tuam: * et justificatiónes tuas doce me.

Servus tuus sum ego: * da mihi intelléctum, ut sciam testimónia tua.

Tempus faciéndi, Dómine: * dissipavérunt legem tuam.

Ideo diléxi mandáta tua, * super aurum et topázion.

Proptérea ad ómnia mandáta tua dirigébar: * omnem viam iníquam ódio hábui.

\rubric{Extra Tempus Paschale: Ant. ut in Proprio de Tempore.}

\variantheading{Tempore Paschali}

\antlabel{Ant.} Allelúja, allelúja, allelúja.

\rubric{Reliqua ut in Proprio de Tempore et in Ordinario. 25}

\sectionornament

\hourheading{Ad Nonam}

Pater noster et Ave.

\V\ Deus, in adjutórium.

\hymnusheading{Hymnus}

\lettrine[lines=2]{R}{ERUM}, Deus, tenax vigor,\\
Immótus in te pérmanens,\\
Lucis diúrnæ témpora\\
Succéssibus detérminans:

Largíre lumen véspere,\\
Quo vita nusquam décidat,\\
Sed præmium mortis sacræ\\
Perénnis instet glória.

* Præsta, Pater piíssime,\\
Patríque compar Unice,\\
Cum Spíritu Paráclito\\
Regnans per omne sæculum. Amen.

\rubric{Extra Tempus Paschale: Ant. ut in Proprio de Tempore.}

\variantheading{Tempore Paschali}

\antlabel{Ant.} Allelúja.

\psalmheading{Psalmus 118, ix}

\lettrine[lines=2]{M}{IRABILIA} testimónia tua: * ídeo scrutáta est ea ánima mea.

Declarátio sermónum tuórum illúminat: * et intelléctum dat párvulis.

Os meum apérui, et attráxi spíritum: * quia mandáta tua desiderábam.

Aspice in me, et miserére mei, * secúndum judícium diligéntium nomen tuum.

Gressus meos dírige secúndum elóquium tuum: * et non dominétur mei omnis injustítia.

Rédime me a calúmniis hóminum: * ut custódiam mandáta tua.

Fáciem tuam illúmina super servum tuum: * et doce me justificatiónes tuas.

Exitus aquárum deduxérunt óculi mei: * quia non custodiérunt legem tuam.

Justus es, Dómine: * et rectum judícium tuum.

Mandásti justítiam testimónia tua: * et veritátem tuam nimis.

Tabéscere me fecit zelus meus: * quia oblíti sunt verba tua inimíci mei.

Ignítum elóquium tuum veheménter: * et servus tuus diléxit illud.

Adolescéntulus sum ego et contémptus: * justificatiónes tuas non sum oblítus.

Justítia tua, justítia in ætérnum: * et lex tua véritas.

Tribulátio, et angústia invenérunt me: * mandáta tua meditátio mea est.

Æquitas testimónia tua in ætérnum: * intelléctum da mihi, et vivam.

\psalmheading{Psalmus 118, x}

\lettrine[lines=2]{C}{LAMAVI} in toto corde meo, exáudi me, Dómine: * justificatiónes tuas requíram.

Clamávi ad te, salvum me fac: * ut custódiam mandáta tua.

Prævéni in maturitáte, et clamávi: * quia in verba tua supersperávi.

Prævenérunt óculi mei ad te dilúculo: * ut meditárer elóquia tua.

Vocem meam audi secúndum misericórdiam tuam, Dómine: * et secúndum judícium tuum vivífica me.

Appropinquavérunt persequéntes me iniquitáti: * a lege autem tua longe facti sunt.

Prope es tu, Dómine: * et omnes viæ tuæ véritas.

Inítio cognóvi de testimóniis tuis: * quia in ætérnum fundásti ea.

Vide humilitátem meam, et éripe me: * quia legem tuam non sum oblítus.

Júdica judícium meum, et rédime me: * propter elóquium tuum vivífica me.

Longe a peccatóribus salus: * quia justificatiónes tuas non exquisiérunt.

Misericórdiæ tuæ multæ, Dómine: * secúndum judícium tuum vivífica me.

Multi qui persequúntur me, et tríbulant me: * a testimóniis tuis non declinávi.

Vidi prævaricántes, et tabescébam: * quia elóquia tua non custodiérunt.

Vide quóniam mandáta tua diléxi, Dómine: * in misericórdia tua vivífica me.

Princípium verbórum tuórum, véritas: * in ætérnum ómnia judícia justítiæ tuæ.

\psalmheading{Psalmus 118, xi}

\lettrine[lines=2]{P}{RINCIPES} persecúti sunt me gratis: * et a verbis tuis formidávit cor meum.

Lætábor ego super elóquia tua: * sicut qui invénit spólia multa.

Iniquitátem ódio hábui, et abominátus sum: * legem autem tuam diléxi.

Sépties in die laudem dixi tibi, * super judícia justítiæ tuæ.

Pax multa diligéntibus legem tuam: * et non est illis scándalum.

Exspectábam salutáre tuum, Dómine: * et mandáta tua diléxi.

Custodívit ánima mea testimónia tua: * et diléxit ea veheménter.

Servávi mandáta tua, et testimónia tua: * quia omnes viæ meæ in conspéctu tuo.

Appropínquet deprecátio mea in conspéctu tuo, Dómine: * juxta elóquium tuum da mihi intelléctum.

Intret postulátio mea in conspéctu tuo: * secúndum elóquium tuum éripe me.

Eructábunt lábia mea hymnum, * cum docúeris me justificatiónes tuas.

Pronuntiábit lingua mea elóquium tuum: * quia ómnia mandáta tua æquitas.

Fiat manus tua ut salvet me: * quóniam mandáta tua elégi.

Concupívi salutáre tuum, Dómine: * et lex tua meditátio mea est.

Vivet ánima mea, et laudábit te: * et judícia tua adjuvábunt me.

Errávi, sicut ovis, quæ périit: * quære servum tuum, quia mandáta tua non sum oblítus.

\rubric{Extra Tempus Paschale: Ant. ut in Proprio de Tempore.}

\variantheading{Tempore Paschali}

\antlabel{Ant.} Allelúja, allelúja, allelúja.

\rubric{Reliqua ut in Proprio de Tempore et in Ordinario. 27}

\sectionornament

% ===== END OF BATCH 4 (Dominica: Ad Primam, Ad Tertiam, Ad Sextam,
% Ad Nonam). Notes: corrected one evident OCR typo, "Miserétur" ->
% "Misereátur" (verna.md line 2846; the other 15 occurrences of this
% formula in verna.md all spell it "Misereátur"). The plain-text
% (non-bold) "Confíteor. Misereátur. Indulgéntiam." at that same line
% is rendered without \formulalabel, since the source gives it no
% bold/label marking there (unlike its 15 other occurrences, all of
% which are at least partly bold) -- preserved as segmented rather
% than normalized to match. Bare trailing numbers "23", "25", "27"
% after "Reliqua ut in Proprio de Tempore et in Ordinario." preserved
% literally, same treatment as "13" in Batch 3. Markdown "---"
% separator lines (verna.md lines 3282, 3423) rendered as
% \sectionornament. Continues with Ad Vesperas =====
